% !TEX program = xelatex
% 理论力学 · 专题讲义 · 从张量到并矢（2026-11-11 独立返工稿）
\documentclass[11pt]{ctexart}
\input{../rule-versions/2.2.2/lecture-deepdive/assets/theme-violet}
\usepackage{tikz}
\usetikzlibrary{arrows.meta}
\setlength{\headheight}{22pt}
\linespread{1.24}
\raggedbottom
\settags{专题 · 张量与并矢}{理论力学}

% 记号：矢量粗斜体，二阶张量粗正体
\newcommand{\ten}[1]{\mathbf{#1}}
\newcommand{\unitd}{\hat{\ten{1}}}
\newcommand{\tr}{\operatorname{tr}}
\newcommand{\figcap}[1]{\par\vspace{2pt}{\small\color{inkgray}#1}\par}
\newcommand{\step}[1]{\par\smallskip\noindent{\color{violetdeep}\bfseries #1}\hspace{0.35em}}
\tikzset{
  ax/.style={inkgray,densely dotted,->},
  vv/.style={thick,->,>={Stealth[length=1.7mm]}},
  vviolet/.style={violet,thick,->,>={Stealth[length=1.7mm]}},
}

\let\orighdr\hdr
\renewcommand{\hdr}[2]{\par\needspace{10\baselineskip}\orighdr{#1}{#2}}
\let\origsubhdr\subhdr
\renewcommand{\subhdr}[1]{\par\needspace{3\baselineskip}\origsubhdr{#1}}

\setlength{\abovedisplayskip}{6pt plus 2pt minus 2pt}
\setlength{\belowdisplayskip}{6pt plus 2pt minus 2pt}
\setlength{\emergencystretch}{1.2em}

\begin{document}

\doctitle{从张量到并矢：九个分量、一台机器、三处力学落点}
{从简单线性代数出发，走到惯量、欧拉方程与应力}
{专题讲义\quad 2026-11-11\quad 主要来源：TM2026-Topic8-P1 p8--p12；RIT Whelan notes 10 的 pp.2--14}

\callout{\textbf{这份材料为谁写、怎么用。}
为“会做简单的矢量分量与矩阵乘法，但还没有系统学过二次型、特征值、张量与并矢”的读者写。
本稿只假设你会：写矢量的分量、做点乘与叉乘、做矩阵乘矢量、取转置；另外知道“导数是一小段时间内改变量的比值极限，积分是把物体切成小块再把贡献加起来”这两句话。其余概念都在第一次用到的地方补，并尽量配一个能口算的小例。
\par\smallskip
\textbf{读法。}标\textbf{（必读）}的节构成主线；标\textbf{（选读）}的段落回答一个额外问题，跳过它不影响后面的阅读，每处都写了“跳过后从哪里继续”。紧跟正文的\textbf{教学示范}给完整演算，用来学方法；\textbf{独立练习}只在正文给题干与考点，详解统一放在附录 A——请先自己做，再翻答案。}

\callout{\textbf{取材与用途声明。}
课件 \texttt{TM2026-Topic8-P1.pdf} p8--p12 给出的是：正交基与求和约定、$\varepsilon_{ijk}$ 与对偶张量 $\vect a^{\times}$，见 p8；张量积（课件原词是 Dyad）、左右作用、转置、把二阶张量读成线性算子 $b=\ten T\vect a$ 或二次型 $E=\vect a\ten T\vect a$、单位张量 $\unitd$，见 p9；旋转算子 $S=\sum_k\vect n'_k\otimes\vect n_k$ 及其性质，见 p10--p11；同一页还用 $S$ 导出了角速度，具体见 p12 的图。
\textbf{惯量张量、欧拉方程、应力张量三节是课件没有展开的延伸落点}：本稿用同一套并矢语言从头推导，惯量部分并与 RIT Whelan notes 10 的 pp.2--14 对照；注意该讲义的对偶张量符号与课件相反，见 §2.6 的注。课件 p13--p18 的欧拉角与运动学例题属于后续的运动学内容，本稿不重复。
\par\smallskip
本稿是在 2026-11-05 版本基础上独立重写的：主题范围（含惯量、欧拉方程、应力与原题）保留，章节顺序、例子、图和表述按本稿重新组织；旧稿只作为范围与题目来源，不作为事实依据。核对中发现的旧稿笔误记录在随附的 \texttt{source-map.json} 与 \texttt{verification.txt}，不在学习正文里展开。}

\hdr{§0}{路线图：三个坎，以及全篇唯一的主线}

\begin{center}
\tabtop
\begin{tabular}{@{}p{0.10\linewidth}p{0.39\linewidth}p{0.43\linewidth}@{}}
\thead{坎} & \thead{你大概会卡在哪} & \thead{本稿在哪里翻过去} \\
\multicolumn{3}{@{}l@{}}{\color{petal}\rule{\linewidth}{0.5pt}}\\
一 & “张量”听起来像新数学，其实只是线性映射的账本 & §0 把矢量、矩阵、换基讲清；§1 讲二次型与主轴；§2 说明张量就是不依赖坐标的那个对象 \\
二 & 二次型、迹、特征值、谱定理没学过，看到“主轴”就停 & §1 从零补：对称部分、椭球、特征值、谱定理、三个不变量 \\
三 & 老师写 $\vect a\vect b$ 或 $\vect a\otimes\vect b$，不知道怎么算，也不知道和力学有什么关系 & §2 给作用规则与三套记号的字典；§3 旋转算子；§4--§6 三个落点；§7 十二道练习 \\
\end{tabular}
\tabbot
\end{center}

\warnbox{\textbf{全篇只有一条主线，先记住它。}
\textbf{同一个物理对象，换坐标系时分量会变，但对象本身不变。}
张量是“对象”，矩阵是“这个对象在一组基下的账本”，并矢是“这个对象的另一种写法”。
后面的每一节都在这句话上做文章：§0 说清账本是什么、换基时怎么变；§1 找出账本里不随基变的数与主轴；§2 用“怎么写、怎么作用”把对象写成并矢；§3--§6 用它处理旋转、惯量、欧拉方程和应力。}

\textbf{最短主线。}如果时间有限，按 §0.3 $\to$ §0.4 $\to$ §1.3 $\to$ §1.5 的定理陈述 $\to$ §2.2 $\to$ §2.3 $\to$ §2.5 $\to$ §2.9 $\to$ §3.2 $\to$ §3.4 $\to$ §4.1 $\to$ §4.3 $\to$ §5.1 $\to$ §5.2 $\to$ §6.1 读，跳过所有标（选读）的段落即可；这些段落里没有后面要用的定义。
若你手上正拿着课件 p9，想马上对上号，可先翻 §2.10 的三记号字典，再回 §2.2 补作用规则。

\hdr{§0}{起点：六个动作与一句话}

这一节不是线性代数课，只把后面反复用的六个动作练一遍：写分量、点乘、叉乘、矩阵乘矢量、换基、算迹与行列式。每个动作配一个能口算的例子。已经熟的读者可以只看每段末尾的结论，但\textbf{建议至少扫一眼 §0.4 的换基}，因为它是全篇的枢纽。

\callout{\textbf{本稿用到的微积分只有两句。}
（1）\textbf{导数}是“改变量除以时间”的极限：$\dd\vect A/\dd t$ 就是 $\vect A$ 在一小段时间里的改变量除以这段时间，再让时间趋于零。
（2）\textbf{积分}是“把物体切成小块，把每块的贡献加起来”：$\int f\dd m$ 读作“把每个质量块的质量 $\dd m$ 乘上它对应的 $f$，再加起来”。均匀物体的 $\dd m=\rho\dd V$，$\rho$ 是密度。
本稿不会用到更复杂的积分技巧；§4.4 只用一次一维积分 $\int_{-L/2}^{L/2}x^2\dd x$，那里会写全。}

\subhdr{0.1\quad 矢量、基与分量（必读）}

取一组互相垂直的单位矢量 $\vect n_1,\vect n_2,\vect n_3$，叫一组\textbf{正交归一基}。任意矢量 $\vect a$ 都能唯一写成
\[
\vect a=a_1\vect n_1+a_2\vect n_2+a_3\vect n_3=\sum_{k=1}^{3}a_k\vect n_k .
\]
三个数 $a_1,a_2,a_3$ 是 $\vect a$ 在这组基下的\textbf{分量}。求法只有一条：拿基矢量去点它，
\[
a_k=\vect a\cdot\vect n_k .
\]
这是因为正交归一基满足 $\vect n_i\cdot\vect n_k=\delta_{ik}$，其中
\[
\delta_{ik}=\begin{cases}1,&i=k,\\0,&i\ne k,\end{cases}
\]
叫\textbf{克罗内克符号}。

\textbf{点乘与投影。}两个矢量的点乘是
\[
\vect a\cdot\vect b=\sum_{k=1}^{3}a_kb_k=a_1b_1+a_2b_2+a_3b_3 .
\]
它是数，不是矢量。若 $\vect n$ 是单位矢量，$\vect b\cdot\vect n$ 就是 $\vect b$ 在 $\vect n$ 方向上的投影长度（带符号）；投影矢量是 $(\vect b\cdot\vect n)\vect n$。

\textbf{叉乘。}两个矢量的叉乘还是矢量，方向按右手定则，分量写成
\[
\vect a\times\vect b=
\begin{pmatrix}
a_2b_3-a_3b_2\\
a_3b_1-a_1b_3\\
a_1b_2-a_2b_1
\end{pmatrix}.
\]
它满足 $\vect a\times\vect b=-\vect b\times\vect a$，且 $\vect a\cdot(\vect a\times\vect b)=0$、$\vect b\cdot(\vect a\times\vect b)=0$（叉乘结果与两个因子都垂直）。§2.6 会用 $\varepsilon_{ijk}$ 把它写得更紧凑，并说明叉乘其实是一个反对称张量在作用；这里先用分量式。

\textbf{求和约定。}为了少写求和号，约定：\textbf{同一项里同一个指标出现两次，就自动对 $1,2,3$ 求和}。例如
\[
\vect a\cdot\vect b=a_kb_k,\qquad
(\ten T\vect x)_i=T_{ij}x_j,\qquad
a_i=\delta_{ij}a_j .
\]
注意只有重复两次才求和；同一个指标出现三次（如 $a_ib_ic_i$ 里 $i$ 出现三次）不是约定用法，遇到时要写清楚对谁求和。本稿在第一次出现时仍常写 $\sum$，读者熟悉后再逐步省略。

\textbf{数字小例。}取 $\vect a=(1,2,2)$、$\vect b=(2,1,0)$。则
\[
|\vect a|=\sqrt{1+4+4}=3,\qquad
\vect a\cdot\vect b=2+2+0=4,
\]
$\vect b$ 在 $\vect a$ 上的投影矢量是 $\dfrac{4}{9}(1,2,2)$；叉乘
\[
\vect a\times\vect b=(2\cdot0-2\cdot1,\;2\cdot2-1\cdot0,\;1\cdot1-2\cdot2)=(-2,4,-3),
\]
核对：$\vect a\cdot(\vect a\times\vect b)=-2+8-6=0$，且 $|\vect a\times\vect b|=\sqrt{4+16+9}=\sqrt{29}$。这三条都可以自己口算一遍。

\subhdr{0.2\quad 线性映射：给一个矢量，吐一个矢量（必读）}

理论力学里到处是“给一个矢量，吐一个矢量”的关系。先看三类常见的量：

\begin{center}
\tabtop
\begin{tabular}{@{}p{0.30\linewidth}p{0.20\linewidth}p{0.40\linewidth}@{}}
\thead{关系} & \thead{进} & \thead{出} \\
\multicolumn{3}{@{}l@{}}{\color{petal}\rule{\linewidth}{0.5pt}}\\
功 $W=\vect F\cdot\vect r$ & 一个矢量 $\vect r$ & 一个数 \\
力矩 $\vect M=\vect r\times\vect F$ & 两个矢量 & 一个矢量 \\
角动量 $\vect L=\ten J\cdot\vect\omega$ & 一个矢量 $\vect\omega$ & 一个矢量 \\
\end{tabular}
\tabbot
\end{center}

第一行“吃矢量吐数”，是线性函数；把“数”看成一维矢量，它的定义与下面的线性映射相同，只是输出是一维。第三行“吃矢量吐矢量”，才是本稿要用的\textbf{二阶张量}。第二行“吃两个矢量吐一个矢量”，可以先固定一个再看成线性映射。

\begin{center}
\fcolorbox{violet}{bandviolet}{\parbox{0.94\linewidth}{\centering\vspace{3pt}
\small 定义：映射 $\ten T$ 若对任意矢量 $\vect u,\vect v$ 与任意数 $\alpha,\beta$ 都满足
\[
\ten T(\alpha\vect u+\beta\vect v)=\alpha\,\ten T(\vect u)+\beta\,\ten T(\vect v),
\]
就称 $\ten T$ 是\textbf{线性映射}。\vspace{3pt}}}
\end{center}

线性映射省事到极点：\textbf{只要知道它作用在三个基矢量上的结果，就知道它对一切矢量的结果。}理由是任意矢量都能写成 $\vect x=\sum_j x_j\vect e_j$，于是由线性
\[
\ten T(\vect x)=\ten T\Bigl(\sum_j x_j\vect e_j\Bigr)=\sum_j x_j\,\ten T(\vect e_j).
\]
右边只用到三个矢量 $\ten T(\vect e_1),\ten T(\vect e_2),\ten T(\vect e_3)$。每个 $\ten T(\vect e_j)$ 本身是一个矢量，要三个数才能写出来，所以\textbf{一共九个数的信息量}。这就是“二阶张量有九个分量”的全部来源：三个输入基矢量，每个吐三个分量。

\subhdr{0.3\quad 矩阵是这台机器的账本（必读）}

把 $\ten T(\vect e_j)$ 在基 $\{\vect e_i\}$ 下展开，系数记作 $T_{ij}$：
\[
\ten T(\vect e_j)=\sum_{i=1}^{3}T_{ij}\vect e_i
\qquad\Longleftrightarrow\qquad
\boxed{\;(\ten T\vect x)_i=\sum_{j=1}^{3}T_{ij}x_j\;}
\]
把九个 $T_{ij}$ 按“第 $i$ 行、第 $j$ 列”排成方阵，就是 $\ten T$ 的矩阵。\textbf{指标的位置是有含义的：第二个指标 $j$ 是“输入基矢量”的编号，第一个指标 $i$ 是“输出分量”的编号。} 换句话说，\textbf{矩阵的第 $j$ 列就是 $\ten T(\vect e_j)$ 的三个分量}。这个约定决定了后面所有点乘的方向，混了就全错。

\textbf{数字例子。}设 $\ten T$ 把 $\vect e_1$ 变成 $2\vect e_1$，把 $\vect e_2$ 变成 $\vect e_1+\vect e_2$，把 $\vect e_3$ 变成 $3\vect e_3$：
\[
\ten T(\vect e_1)=\begin{pmatrix}2\\0\\0\end{pmatrix},\quad
\ten T(\vect e_2)=\begin{pmatrix}1\\1\\0\end{pmatrix},\quad
\ten T(\vect e_3)=\begin{pmatrix}0\\0\\3\end{pmatrix}
\quad\Longrightarrow\quad
\ten T=\begin{pmatrix}2&1&0\\0&1&0\\0&0&3\end{pmatrix}.
\]
注意这三个列矢量是按列摆进去的。取 $\vect x=(1,2,3)$：
\[
\ten T\vect x=\begin{pmatrix}2&1&0\\0&1&0\\0&0&3\end{pmatrix}\begin{pmatrix}1\\2\\3\end{pmatrix}
=\begin{pmatrix}2\cdot1+1\cdot2\\0\cdot1+1\cdot2\\0+0+3\cdot3\end{pmatrix}=\begin{pmatrix}4\\2\\9\end{pmatrix}.
\]
口算核对：$\ten T(\vect x)=1\cdot(2,0,0)+2\cdot(1,1,0)+3\cdot(0,0,3)=(4,2,9)$。两条路一致。

\textbf{矩阵乘法就是复合。}若先作用 $\ten B$ 再作用 $\ten A$，则 $(\ten A\ten B)\vect x=\ten A(\ten B\vect x)$，分量式是
\[
(\ten A\ten B)_{ij}=\sum_k A_{ik}B_{kj}.
\]
\textbf{转置}是行列互换：$(\ten A^{\mathsf T})_{ij}=A_{ji}$。它有一条后面反复用的判别式：
\[
(\ten A\vect x)\cdot\vect y=\vect x\cdot(\ten A^{\mathsf T}\vect y).
\]
证明只有一行：左边是 $A_{ij}x_jy_i$，右边是 $x_jA_{ij}y_i$，只是把数相乘的顺序换了。

\subhdr{0.4\quad 换基：账本变，机器不变（必读）}

同一台机器，换一组基来记账，九个数字会变。这句话是本稿的枢纽，值得用一个具体例子钉死。

\textbf{二维例子。}机器 $\ten T$ 是“沿 $x$ 轴拉伸两倍”：$\ten T(x,y)=(2x,y)$。在标准基 $\vect e_1,\vect e_2$ 下
\[
\ten T=\begin{pmatrix}2&0\\0&1\end{pmatrix}.
\]
现在改用转了 $45^\circ$ 的基 $\vect n_1=(\vect e_1+\vect e_2)/\sqrt2$、$\vect n_2=(-\vect e_1+\vect e_2)/\sqrt2$。
先算输入：$\ten T\vect n_1=\ten T\bigl(\tfrac{1}{\sqrt2},\tfrac{1}{\sqrt2}\bigr)=\bigl(\tfrac{2}{\sqrt2},\tfrac{1}{\sqrt2}\bigr)$。
再把它在这组新基下分解，用 $\vect n_i\cdot$ 投影取系数：
\[
\ten T\vect n_1\cdot\vect n_1=\frac{2+1}{2}=\frac32,\qquad
\ten T\vect n_1\cdot\vect n_2=\frac{-2+1}{2}=-\frac12 .
\]
同法 $\ten T\vect n_2=\bigl(-\tfrac{2}{\sqrt2},\tfrac{1}{\sqrt2}\bigr)$，
\[
\ten T\vect n_2\cdot\vect n_1=\frac{-2+1}{2}=-\frac12,\qquad
\ten T\vect n_2\cdot\vect n_2=\frac{2+1}{2}=\frac32 .
\]
于是同一台机器的账本变成
\[
\ten T'=\begin{pmatrix}3/2&-1/2\\-1/2&3/2\end{pmatrix}.
\]
两个账本长得完全不一样，可它们描述的是同一件事：沿 $\vect e_1$ 方向拉两倍、沿 $\vect e_2$ 方向不变。\textbf{判断“是不是同一台机器”，不能看数字，要看数字怎么变。}下面把这条换算规矩推出来，只有三步。

\step{第一步：每个新分量都是两次点乘}
回看上面的算法，新账本的每一个元素都长成同一个形状：
\[
T'_{ij}=\vect n_i\cdot(\ten T\vect n_j).
\]
内层那个点乘在\textbf{输入侧}：$\ten T$ 只认旧基坐标，所以得先把 $\vect n_j$ 写成旧基里的样子；外层那个点乘在\textbf{输出侧}：$\ten T$ 吐出来的仍是旧基坐标，得再投影到 $\vect n_i$ 上，才算新基下的第 $i$ 个分量。
\textbf{两次翻译，就是后面那两个求和。}向量换基只乘一次，映射换基要乘两次：向量只有“分量”一个位置要翻，映射有“输入”和“输出”两个位置，各翻一次。

\step{第二步：把两次点乘按旧基写开}
记
\[
R_{ki}=\vect n_i\cdot\vect e_k,
\]
它的第 $i$ 列是 $\vect n_i$ 在旧基下的三个分量。于是 $\vect n_j=\sum_l R_{lj}\vect e_l$，喂给 $\ten T$ 之后
\[
\ten T\vect n_j=\sum_k\Bigl(\sum_l T_{kl}R_{lj}\Bigr)\vect e_k .
\]
最后拿 $\vect n_i=\sum_k R_{ki}\vect e_k$ 去点它，$R_{ki}$ 就是第二个因子：
\[
T'_{ij}=\sum_{k}R_{ki}\Bigl(\sum_{l}T_{kl}R_{lj}\Bigr)=\sum_{k,l}R_{ki}T_{kl}R_{lj}.
\]

\step{第三步：读成两个矩阵乘}
\[
\boxed{\;\ten T'=\ten R^{\mathsf T}\ten T\ten R\;},\qquad R_{ki}=\vect n_i\cdot\vect e_k .
\]
$R$ 的第 $i$ 列是 $\vect n_i$ 在旧基下的分量；正交基之间 $R^{\mathsf T}R=\unitd$（下面 §3.3 会再证一次），所以 $R^{-1}=R^{\mathsf T}$。这个等式就是“新基也正交归一”的分量写法：$R^{\mathsf T}R$ 的 $(i,k)$ 元是 $\vect n_i\cdot\vect n_k=\delta_{ik}$。这个夹心结构怎么读，只看每个因子贴着哪一侧：

\begin{center}
\tabtop
\begin{tabular}{@{}p{0.24\linewidth}p{0.20\linewidth}p{0.48\linewidth}@{}}
\thead{因子} & \thead{贴在} & \thead{它在干什么} \\
\multicolumn{3}{@{}l@{}}{\color{petal}\rule{\linewidth}{0.5pt}}\\
右边的 $\ten R$ & 输入侧 & 把新基坐标翻成旧基坐标，$\ten T$ 才认 \\
中间的 $\ten T$ & 机器本身 & 它只会在旧账本上算 \\
左边的 $\ten R^{\mathsf T}$ & 输出侧 & 算完再翻回新基坐标 \\
\end{tabular}
\tabbot
\end{center}

一句话口诀：\textbf{指标在哪一侧，翻译就贴在那一侧}——$i$ 是输出指标，配左边的 $\ten R^{\mathsf T}$；$j$ 是输入指标，配右边的 $\ten R$。
顺序不能对调：$\ten R\ten T\ten R^{\mathsf T}$ 是\textbf{把基向相反方向转同样的角度}后得到的账本，它与 $\ten T'$ 相似（迹、行列式、特征值都一样），但不是同一组基下的账本。

\textbf{矢量的新坐标是 $R^{\mathsf T}$ 乘旧坐标。}这一点最容易写反，单独说清：
\[
x'_i=\vect n_i\cdot\vect x=\sum_k(\vect n_i\cdot\vect e_k)x_k=\sum_k R_{ki}x_k=(R^{\mathsf T}\vect x)_i .
\]
所以\textbf{矢量贴一个 $R^{\mathsf T}$，二阶张量左右各贴一个}。如果你读的书把 $R$ 定义成“行是新基分量”，公式会写成 $\ten T'=\ten R\ten T\ten R^{\mathsf T}$、$\vect x'=\ten R\vect x$；两套只差一个转置，不要混用。本稿固定用“列是新基分量”的 $R$。

\textbf{核对。}上面例子里 $R=\dfrac{1}{\sqrt2}\begin{pmatrix}1&-1\\1&1\end{pmatrix}$，代入 $\ten R^{\mathsf T}\ten T\ten R$ 得 $\begin{pmatrix}3/2&-1/2\\-1/2&3/2\end{pmatrix}$，与逐项投影所得一致。
\aside{这只是\textbf{核对}：规矩本身由上面三步推出，不是靠这个数值例子“试出来”的。§1.5 的谱定理会再给一个更强的结论：对对称 $\ten T$，总可以选一组基使 $\ten T'$ 变成对角阵；那组基就是主轴。}

\subhdr{0.5\quad 账本里不随基变的数：迹、行列式、特征多项式（必读）}

既然换个基数字全变，一个自然的问题是：\textbf{账本里有没有哪个组合换基后不变？}有，而且恰好是需要的那几个。

\textbf{迹。}方阵 $\ten A$ 的\textbf{迹}是对角元之和：
\[
\tr\ten A=\sum_i A_{ii}.
\]
它换基不变，理由只有一行。矩阵乘法有“循环性”：
\[
\tr(\ten A\ten B)=\sum_i\sum_j A_{ij}B_{ji}=\sum_j\sum_i B_{ji}A_{ij}=\tr(\ten B\ten A).
\]
取 $\ten B=\ten R$、$\ten A=\ten R^{\mathsf T}\ten T$，就有
\[
\tr(\ten R^{\mathsf T}\ten T\ten R)
=\tr\bigl((\ten R^{\mathsf T}\ten T)\ten R\bigr)
=\tr\bigl(\ten R(\ten R^{\mathsf T}\ten T)\bigr)
=\tr\bigl((\ten R\ten R^{\mathsf T})\ten T\bigr)
=\tr\ten T .
\]
最后一步用了 $\ten R\ten R^{\mathsf T}=\unitd$（正交基的性质）。
\textbf{一行读完：迹不变，因为迹允许把因子转圈，而正交矩阵转一圈就消掉了。}

\textbf{数字核对。}§0.4 的两个账本：
\[
\tr\begin{pmatrix}2&0\\0&1\end{pmatrix}=2+1=3,
\qquad
\tr\begin{pmatrix}3/2&-1/2\\-1/2&3/2\end{pmatrix}=\frac32+\frac32=3 .
\]

\textbf{行列式。}$\det\ten T$ 是这台机器对\textbf{体积}的缩放倍数：把单位立方体喂给 $\ten T$，出来一个平行六面体，它的体积是 $|\det\ten T|$；符号表示“是否翻转了手性”（手性就是左右手的区别，符号为负表示像照镜子那样翻转）。这个图像立刻给出两件事：
\begin{itemize}
\item 换坐标系不会改变“体积被放大几倍”，所以 $\det$ 也是不变量。代数上同样一句话：
\[
\det(\ten R^{\mathsf T}\ten T\ten R)=\det(\ten R^{\mathsf T})\det\ten T\det(\ten R)=\det\ten T,
\]
因为正交矩阵满足 $\det\ten R=\pm1$，平方为 $1$。
\item $\det\ten T=0$ 意味着体积被压成零，机器不可逆；后面“某个主值为零”就是这种退化。
\end{itemize}
二维行列式的算法是 $\det\begin{pmatrix}a&b\\c&d\end{pmatrix}=ad-bc$。三维按第一行展开：
\[
\det\begin{pmatrix}a&b&c\\d&e&f\\g&h&i\end{pmatrix}
=a(ei-fh)-b(di-fg)+c(dh-eg).
\]
例如 $\det\begin{pmatrix}4&1&0\\1&4&0\\0&0&2\end{pmatrix}
=4(8-0)-1(2-0)+0=30$，这个数在 P1 里会再用到。本稿只在具体数字里用到三维行列式，不要求背展开式，只要会按一行“正、负、正”交替地展开即可。

\textbf{数字核对。}§0.4 里 $\ten T=\operatorname{diag}(2,1)$，$\det\ten T=2$：面积放大两倍。换基后
\[
\ten T'=\begin{pmatrix}3/2&-1/2\\-1/2&3/2\end{pmatrix},
\qquad
\det\ten T'=\frac94-\frac14=2 .
\]

\textbf{特征多项式。}把“$\det$ 与 $\tr$ 都不变”这个念头推到底，就是问：$\det(\ten T-\lambda\unitd)$ 作为 $\lambda$ 的多项式是什么？它的系数是不变量，根就是特征值（§1.4 会正式使用）。
二维先看清楚。$\ten T=\begin{pmatrix}a&b\\c&d\end{pmatrix}$，
\[
\det(\ten T-\lambda\unitd)=\det\begin{pmatrix}a-\lambda&b\\c&d-\lambda\end{pmatrix}
=(a-\lambda)(d-\lambda)-bc=\lambda^2-(a+d)\lambda+(ad-bc).
\]
所以二维的式子只有两个系数，分别就是 $\tr$ 与 $\det$：
\[
\lambda^2-(\tr\ten T)\lambda+\det\ten T=0 .
\]
三维照做，得到一个三次多项式：
\[
\det(\ten T-\lambda\unitd)=-\lambda^3+I_1\lambda^2-I_2\lambda+I_3 ,
\]
\[
I_1=\tr\ten T,\qquad
I_3=\det\ten T,\qquad
I_2=\tfrac12\bigl[(\tr\ten T)^2-\tr\ten T^2\bigr].
\]
三个系数 $I_1,I_2,I_3$ 都由 $\tr$ 与 $\det$ 拼成，所以\textbf{三个都不随基变}。
\aside{为什么要有 $\ten T^2$：看二维就明白。$\tr\ten T^2=a^2+2bc+d^2$，$(\tr\ten T)^2=a^2+2ad+d^2$，相减除二得 $ad-bc=\det\ten T$。所以 $I_2$ 不是新东西，是“二维里 $\det$ 的位置在三维里的替身”。三维有三个独立的不变量，需要三个组合。}

\textbf{数字小例。}取 $\ten T=\operatorname{diag}(1,2,3)$。则 $I_1=6$、$I_3=6$、$\tr\ten T^2=1+4+9=14$、$I_2=\tfrac12(36-14)=11$。特征多项式是
\[
-\lambda^3+6\lambda^2-11\lambda+6=0,
\]
即 $(\lambda-1)(\lambda-2)(\lambda-3)=0$，根正好是 $1,2,3$；系数也正好是“三数之和 $6$、两两乘积之和 $11$、三数乘积 $6$”。这三个组合在换基后仍然是它们。

\warnbox{\textbf{一个容易误解的地方。}
“$I_1,I_2,I_3$ 不变”说的是\textbf{换基时}不变，不是“$\ten T$ 不变”。换基时九个分量全变，只有这三个组合不变。特征值与主轴方向的关系在 §1.4--§1.6 讲；那里会说明：特征值不变，但主轴方向的\textbf{分量}会随基一起转。}

\hdr{§1}{二次型与主轴：椭球、特征值、谱定理}

这一节把“二次型—椭圆—特征值—主轴”连成一条线。目标不是会证明谱定理，而是会用它的三条结论（特征值与特征向量的正式定义见 §1.4）：特征值都是实数、不同特征值的特征向量互相垂直、存在一组正交主轴使账本变成对角阵。会用这三条，后面的惯量与应力就只是代入。

\subhdr{1.1\quad 二阶张量的两个身份：算子与二次型（必读）}

同一个二阶张量有两种读法，课件 p9 把两种都写了：
\[
\text{作为线性算子：}\quad \vect b=\ten T\vect a,\qquad b_i=T_{ij}a_j ;
\]
\[
\text{作为二次型：}\quad E=\vect a\cdot\ten T\vect a=T_{ik}a_ia_k .
\]
第一种读法“吃一个矢量，吐一个矢量”，§0 已经讲过；第二种读法“吃一个矢量，吐一个数”，本小节讲。

\textbf{“二次的”是什么意思。}把 $\vect a$ 换成 $2\vect a$，$E$ 变成 $4E$；换成 $-\vect a$，$E$ 不变。所以它刻画的是“方向 + 大小”，不是一个新的矢量。三维写开就是
\[
E=\sum_{i=1}^{3}\sum_{k=1}^{3}T_{ik}a_ia_k
=T_{11}a_1^2+T_{22}a_2^2+T_{33}a_3^2
+(T_{12}+T_{21})a_1a_2+(T_{13}+T_{31})a_1a_3+(T_{23}+T_{32})a_2a_3 .
\]
注意交叉项的系数是 $T_{12}+T_{21}$，不是无条件的一个 $2T_{12}$；只有当 $\ten T$ 对称（$T_{12}=T_{21}$）时，才可以写成 $2T_{12}a_1a_2$。这个区别在 §1.2 会变成一条重要结论。

\textbf{物理上它长什么样。}刚体的转动动能
\[
K=\tfrac12\vect\omega\cdot\ten J\cdot\vect\omega=\tfrac12J_{ik}\omega_i\omega_k
\]
就是二次型（差一个 $\tfrac12$）。这句话解释了为什么“惯量”只能用张量描述：能量是 $\vect\omega$ 的二次函数，而二次函数的系数天然是一张九元素的表。

\subhdr{1.2\quad 只有对称部分起作用（必读）}

\textbf{一条必须先钉死的事实：}二次型对 $\ten T$ 的\textbf{反对称部分完全无感}。也就是说，算 $E$ 时可以把 $\ten T$ 换成它自己的对称部分
\[
\ten S=\tfrac12(\ten T+\ten T^{\mathsf T}),\qquad S_{ik}=\tfrac12(T_{ik}+T_{ki}),
\]
结果一字不差。

理由是纯粹的指标改名。交叉项 $\sum_{i,k}T_{ik}a_ia_k$ 里，把 $i$ 与 $k$ 的名字对换（只是换字母，值不变）：
\[
\sum_{i,k}T_{ik}a_ia_k=\sum_{i,k}T_{ki}a_ka_i=\sum_{i,k}T_{ki}a_ia_k .
\]
两式相加再除二：
\[
\sum_{i,k}T_{ik}a_ia_k
=\sum_{i,k}\tfrac12(T_{ik}+T_{ki})a_ia_k
=\sum_{i,k}S_{ik}a_ia_k .
\]
\textbf{数字例子。}取 $\ten T=\begin{pmatrix}1&2\\0&1\end{pmatrix}$（不对称）。算 $E$：
\[
E=\begin{pmatrix}x&y\end{pmatrix}\begin{pmatrix}1&2\\0&1\end{pmatrix}\begin{pmatrix}x\\y\end{pmatrix}
=x^2+2xy+y^2=(x+y)^2 .
\]
换成对称部分 $\tfrac12(\ten T+\ten T^{\mathsf T})=\begin{pmatrix}1&1\\1&1\end{pmatrix}$：
\[
E=\begin{pmatrix}x&y\end{pmatrix}\begin{pmatrix}1&1\\1&1\end{pmatrix}\begin{pmatrix}x\\y\end{pmatrix}=x^2+xy+yx+y^2=x^2+2xy+y^2 .
\]
一样。\textbf{所以从这一节起，凡讲二次型都默认 $\ten T$ 对称。}力学里的 $\ten J$ 与 $\ten\sigma$ 本来就对称（§4.2、§6.2 会各给一次理由），不会损失任何东西。

\subhdr{1.3\quad 二次型的等值面是椭球：主轴就是它的轴（必读）}

令 $E=1$，得到一个曲面。这个曲面的形状决定了一切。先看没有交叉项的情形。

\aside{本节会先借用两个词：若存在非零矢量 $\vect n$ 与数 $\lambda$ 使 $\ten T\vect n=\lambda\vect n$，就把 $\vect n$ 叫\textbf{主轴方向}（正式名称是特征向量），把 $\lambda$ 叫对应的\textbf{主值}（正式名称是特征值）。§1.4 会给出正式定义与求法；这里先用这个含义读椭圆的两个轴。}

\textbf{没有交叉项。}$\ten T=\begin{pmatrix}1&0\\0&4\end{pmatrix}$，则 $E=x^2+4y^2=1$，这是椭圆
\[
\frac{x^2}{1^2}+\frac{y^2}{(1/2)^2}=1,
\]
半轴分别是 $1$ 与 $1/2$。注意半轴是 $1/\sqrt{\lambda}$，$\lambda$ 是 $\ten T$ 的对角元（也是它的特征值）：$\lambda=1$ 给半轴 $1$，$\lambda=4$ 给半轴 $1/2$。\textbf{特征值越大，那个方向越“扁”。}

\textbf{有交叉项。}$\ten T=\begin{pmatrix}2&1\\1&2\end{pmatrix}$，则 $E=2x^2+2xy+2y^2=1$。这个椭圆是斜的。它的两个特征向量是
\[
\vect n_1=\frac{1}{\sqrt2}\begin{pmatrix}1\\1\end{pmatrix}\quad(\lambda=3),
\qquad
\vect n_2=\frac{1}{\sqrt2}\begin{pmatrix}-1\\1\end{pmatrix}\quad(\lambda=1).
\]
改用沿这两个方向的坐标 $u=\vect n_1\cdot(x,y)=\dfrac{x+y}{\sqrt2}$、$v=\vect n_2\cdot(x,y)=\dfrac{-x+y}{\sqrt2}$，反解是
\[
x=\frac{u-v}{\sqrt2},\qquad y=\frac{u+v}{\sqrt2}.
\]
代入：
\[
x^2+y^2=u^2+v^2,\qquad xy=\frac{u^2-v^2}{2},
\]
\[
E=2(u^2+v^2)+2\cdot\frac{u^2-v^2}{2}=3u^2+v^2 .
\]
\textbf{交叉项消失了}，曲面变成半轴 $1/\sqrt3$ 与 $1$ 的正椭圆：沿 $u$（即 $\vect n_1$）的半轴是 $1/\sqrt3$，沿 $v$（即 $\vect n_2$）的半轴是 $1$。让交叉项消失的那组坐标轴，就是\textbf{主轴}。

\begin{center}
\begin{tikzpicture}[scale=2.15,>={Stealth[length=1.7mm]}]
  \draw[inkgray,densely dotted,->] (-1.25,0)--(1.3,0) node[right]{\scriptsize $x$};
  \draw[inkgray,densely dotted,->] (0,-1.2)--(0,1.3) node[above]{\scriptsize $y$};
  \draw[violet,thick,fill=bandviolet,rotate around={45:(0,0)}] (0,0) ellipse ({0.5774} and {1});
  \draw[violet,thick,->] (0,0)--(0.85,0.85) node[above right]{\scriptsize $u$（$\vect n_1$）};
  \draw[violet,thick,->] (0,0)--(-0.85,0.85) node[above left]{\scriptsize $v$（$\vect n_2$）};
  \node[violet] at (0.60,0.33) {\scriptsize $\lambda=3$};
  \node[violet] at (-0.62,0.33) {\scriptsize $\lambda=1$};
  \fill (0,0) circle (0.7pt);
\end{tikzpicture}\\[2pt]
\figcap{\mbox{\textbf{图 1}（自绘）：}$E=2x^2+2xy+2y^2=1$ 的等值线。灰虚线是原来的 $x,y$ 轴（在这组轴上二次型带交叉项，椭圆是“斜”的）；紫线是主轴 $u=\vect n_1$、$v=\vect n_2$。沿 $u$ 的特征值是 $3$，半轴只有 $1/\sqrt3\approx0.577$；沿 $v$ 的特征值是 $1$，半轴是 $1$。一句话：\textbf{半轴 $=1/\sqrt{\lambda}$，特征值越大那一轴越短}。}
\end{center}

这不是巧合，而是全部内容：\textbf{对任何对称 $\ten T$，总存在一组正交坐标轴，使二次型没有交叉项；在那一组轴上，系数就是特征值。}这句话的严格版本就是 §1.5 的谱定理。图中 $u$ 沿 $(1,1)/\sqrt2$、$v$ 沿 $(-1,1)/\sqrt2$，与上面代入的 $x=(u-v)/\sqrt2,\;y=(u+v)/\sqrt2$ 逐项对应：$v=0$ 时 $(x,y)=u(1,1)/\sqrt2$ 是 $\vect n_1$ 方向，$u=0$ 时 $(x,y)=v(-1,1)/\sqrt2$ 是 $\vect n_2$ 方向。

\aside{若某个特征值为负，$E=1$ 的等值面不再是椭圆，而是双曲线或空集；若某个特征值为零，等值面退化（例如变成一对平行线或柱面）。但“主轴 = 特征向量”仍然成立。惯量张量通常正定（定义见 §4.2），所以惯量椭球通常是真椭球；退化情形（例如全部质量落在一条过原点的直线上，或细杆沿自身方向）会有一个主值为零，椭球退化成柱面。这里要分清：使用 §1.5 的谱定理只要求 $\ten T$ \textbf{实对称}，并不要求正定；正定只保证 $E=1$ 是椭球。应力张量可以有负的主值，那时用主轴与主应力，不画椭球。}

\subhdr{1.4\quad 特征值与特征向量：主轴方向的定义（必读）}

\begin{center}
\fcolorbox{violet}{bandviolet}{\parbox{0.94\linewidth}{\centering\vspace{3pt}
\small 若存在非零矢量 $\vect n$ 与数 $\lambda$ 使 $\ten T\vect n=\lambda\vect n$，则称 $\lambda$ 为\textbf{特征值}、$\vect n$ 为对应的\textbf{特征向量}。\vspace{3pt}}}
\end{center}

几何含义：\textbf{$\ten T$ 作用在 $\vect n$ 上不改变方向，只改变长度。}这就是“主轴”三个字的全部意思——沿这根轴，机器只做伸缩。

怎么求？把 $\ten T\vect n=\lambda\vect n$ 移项成 $(\ten T-\lambda\unitd)\vect n=\vect0$。这是齐次线性方程组，要它有\textbf{非零}解，系数矩阵必须不可逆（否则唯一解是 $\vect n=\vect0$），也就是
\[
\det(\ten T-\lambda\unitd)=0 .
\]
这正是 §0.5 那个多项式等于零。所以求特征值就是解一个三次方程；\textbf{而它的三个系数正是三个不变量}。

\textbf{三维教学示范。}取
\[
\ten D=\begin{pmatrix}6&2&0\\2&6&0\\0&0&1\end{pmatrix}.
\]
先算特征多项式：
\[
\det(\ten D-\lambda\unitd)
=(1-\lambda)\bigl[(6-\lambda)^2-4\bigr]
=(1-\lambda)(\lambda-4)(\lambda-8).
\]
所以特征值是 $\lambda_1=8,\lambda_2=4,\lambda_3=1$。再求方向：
\begin{itemize}
\item $\lambda=8$：$(\ten D-8\unitd)\vect n=\vect0$ 的前两行给出 $-2x+2y=0$，即 $x=y$，第三行给 $-7z=0$，所以 $\vect n_1=(1,1,0)/\sqrt2$。
\item $\lambda=4$：前两行给出 $2x+2y=0$，即 $y=-x$，第三行给 $3z=0$，所以 $\vect n_2=(1,-1,0)/\sqrt2$。
\item $\lambda=1$：前两行给出 $5x+2y=0$ 与 $2x+5y=0$，只能 $x=y=0$，所以 $\vect n_3=\vect e_3$。
\end{itemize}
三个方向互相垂直、都是单位矢量。把它们按列乘行的方式拼回去：
\[
8\,\vect n_1\vect n_1^{\mathsf T}+4\,\vect n_2\vect n_2^{\mathsf T}+1\,\vect e_3\vect e_3^{\mathsf T}
=8\cdot\frac12\begin{pmatrix}1&1&0\\1&1&0\\0&0&0\end{pmatrix}
+4\cdot\frac12\begin{pmatrix}1&-1&0\\-1&1&0\\0&0&0\end{pmatrix}
+\begin{pmatrix}0&0&0\\0&0&0\\0&0&1\end{pmatrix}
=\begin{pmatrix}6&2&0\\2&6&0\\0&0&1\end{pmatrix}.
\]
这里 $\vect n\vect n^{\mathsf T}$ 是“列乘行”得到的 $3\times3$ 矩阵，$(i,j)$ 元是 $n_in_j$；§2.5 会说明它就是“投影到 $\vect n$ 方向”的并矢。特征向量的正负号可以随便取，取反仍是同一个方向。

\textbf{特征值重复时怎么办。}例如 $\ten D=\operatorname{diag}(2,2,5)$：$\lambda=2$ 对应整个 $xy$ 平面，平面里任意一对正交单位矢量都可以当主轴；$\lambda=5$ 对应 $z$ 轴。重复特征值不破坏谱定理，只让“主轴的选择”不唯一。§4.3 的圆柱、§5.3 的对称陀螺都会用到这一点。

\subhdr{1.5\quad 谱定理：完整陈述与用法（必读）}

\begin{center}
\fcolorbox{violet}{bandviolet}{\parbox{0.94\linewidth}{\vspace{3pt}
\textbf{谱定理（实对称矩阵，三维）。}设 $\ten T$ 是 $3\times3$ 实矩阵且 $\ten T=\ten T^{\mathsf T}$。则
\begin{enumerate}[label=(\roman*),leftmargin=1.8em,itemsep=2pt]
\item $\ten T$ 的三个特征值 $\lambda_1,\lambda_2,\lambda_3$ 全是实数；
\item 属于不同特征值的特征向量互相正交；
\item 存在三个互相正交的单位矢量 $\vect n_1,\vect n_2,\vect n_3$（三个主轴），使
\[
\ten T=\lambda_1\vect n_1\vect n_1^{\mathsf T}+\lambda_2\vect n_2\vect n_2^{\mathsf T}+\lambda_3\vect n_3\vect n_3^{\mathsf T},
\]
即：在以 $\{\vect n_k\}$ 为基的坐标系里，$\ten T$ 的账本就是 $\operatorname{diag}(\lambda_1,\lambda_2,\lambda_3)$。
\end{enumerate}
\vspace{3pt}}}
\end{center}

这里的 $\vect n_k\vect n_k^{\mathsf T}$ 是列乘行矩阵；§2.5 会把它写成并矢 $\vect n_k\otimes\vect n_k$，并说明它是“投影到 $\vect n_k$”的算子。现在只需知道：三项相加就是上面 §1.4 做过的那种拼法。

\textbf{(ii) 的证明只有两行。}设 $\ten T\vect u=\lambda\vect u$，$\ten T\vect v=\mu\vect v$，$\lambda\ne\mu$。则
\[
\lambda(\vect u\cdot\vect v)=(\ten T\vect u)\cdot\vect v=\vect u\cdot(\ten T^{\mathsf T}\vect v)=\vect u\cdot(\ten T\vect v)=\mu(\vect u\cdot\vect v).
\]
第二个等号用了转置判别式 $(\ten T\vect u)\cdot\vect v=\vect u\cdot(\ten T^{\mathsf T}\vect v)$（§0.3），第三个用了对称 $\ten T^{\mathsf T}=\ten T$。于是 $(\lambda-\mu)(\vect u\cdot\vect v)=0$，而 $\lambda\ne\mu$，只能 $\vect u\cdot\vect v=0$。

\textbf{怎么用 (iii)。}把 (iii) 代进二次型，得到一条很好用的公式：
\[
E=\vect a\cdot\ten T\vect a=\lambda_1(\vect n_1\cdot\vect a)^2+\lambda_2(\vect n_2\cdot\vect a)^2+\lambda_3(\vect n_3\cdot\vect a)^2 .
\]
推导只有一步：$\vect a=\sum_k(\vect n_k\cdot\vect a)\vect n_k$，而 $\vect n_k\vect n_k^{\mathsf T}\vect a=\vect n_k(\vect n_k\cdot\vect a)$。这条公式把“二次型”变成“三个投影长度的平方，各乘一个主值”，§4 算绕任意方向的转动惯量时直接用它。

\textbf{瑞利商（推论）。}对任意非零矢量 $\vect a$，
\[
\lambda_{\min}\le\frac{\vect a\cdot\ten T\vect a}{\vect a\cdot\vect a}\le\lambda_{\max}.
\]
这是因为右边是一个以 $(\vect n_k\cdot\vect a)^2/|\vect a|^2$ 为权重的加权平均，权重非负且和为 $1$。§4.3 用它说明“绕任意方向的转动惯量介于最小与最大主值之间”。

\warnbox{\textbf{选读（可跳过）：谱定理 (i)+(iii) 的证明思路。}
这一段回答“为什么一定存在三个正交主轴”，不影响后面使用；跳过后从 §1.6 继续。
\par\smallskip
考虑在单位球面 $\vect x\cdot\vect x=1$ 上求二次型 $f(\vect x)=\vect x\cdot\ten T\vect x$ 的最大值。球面是闭且有界的，$f$ 连续，所以最大值一定取到，记在 $\vect n_1$ 处。用拉格朗日乘子法（约束极值的条件是两个梯度平行）：要 $\nabla f=\mu\nabla g$，其中 $g=\vect x\cdot\vect x-1$。而
\[
\nabla f=2\ten T\vect x,\qquad \nabla g=2\vect x
\quad\Longrightarrow\quad
\ten T\vect n_1=\mu\vect n_1 .
\]
所以极值点必是特征向量，乘子 $\mu$ 就是特征值；并且在该点 $f=\vect n_1\cdot\ten T\vect n_1=\mu\,|\vect n_1|^2=\mu$。这里 $\mu$ 是实数，因为它是实函数在实点上的取值——这就是 (i) 的第一步。
\par\smallskip
接着在“与 $\vect n_1$ 正交的那个平面”里重复同样的论证。需要先确认 $\ten T$ 把这个平面映到自身：若 $\vect x\cdot\vect n_1=0$，则
\[
(\ten T\vect x)\cdot\vect n_1=\vect x\cdot(\ten T\vect n_1)=\mu\,\vect x\cdot\vect n_1=0,
\]
所以 $\ten T\vect x$ 仍在这个平面里。平面是二维的，同样的“闭有界 + 连续”论证给出 $\vect n_2$；再降一次维得到 $\vect n_3$。三个方向两两正交，单位化后就是主轴。
\par\smallskip
最后验证 (iii) 的分解式：两边作用在基矢量 $\vect n_j$ 上，左边给 $\lambda_j\vect n_j$，右边给 $\lambda_j\vect n_j$（因为 $\vect n_k\vect n_k^{\mathsf T}\vect n_j=\delta_{kj}\vect n_k$）。两个矩阵在三个基矢量上作用相同，就相等。证完。}

\subhdr{1.6\quad 三个不变量：验算的最好工具（必读）}

由 §0.5，三个系数是
\[
I_1=\tr\ten T=\lambda_1+\lambda_2+\lambda_3,
\qquad
I_3=\det\ten T=\lambda_1\lambda_2\lambda_3,
\]
\[
I_2=\tfrac12\bigl[(\tr\ten T)^2-\tr\ten T^2\bigr]
=\lambda_1\lambda_2+\lambda_2\lambda_3+\lambda_3\lambda_1 .
\]
右边那两个用 $\lambda$ 表示的等号不需要另证：在主轴系里 $\ten T=\operatorname{diag}(\lambda_1,\lambda_2,\lambda_3)$，而 $\tr$ 与 $\det$ 都不随基变（§0.5），直接代进去就是。

\textbf{用法：做完一道特征值题，用这三个数各查一遍。}以 §1.4 的 $\ten D$ 为例：$\tr\ten D=13$、$\det\ten D=32$、$\tr\ten D^2=81$，所以 $I_2=\tfrac12(169-81)=44$。而算出的特征值是 $8,4,1$：
\[
8+4+1=13,\qquad 8\cdot4+4\cdot1+1\cdot8=44,\qquad 8\cdot4\cdot1=32,
\]
三样都对上。这个“和、两两和、积”的检查几乎不花时间，却能把大部分符号错误当场抓住。

\warnbox{\textbf{三个不变量不是“随便三个数”。}它们分别是特征值的和、两两乘积之和、乘积；换基时它们不变，但\textbf{单个矩阵元会变}。另外，特征值不变而特征向量随基转——说“主轴方向”时必须连着坐标系（或连着刚体）说。}

\hdr{§2}{张量与并矢：对象、写法、运算}

这一节把二阶张量写成并矢，并建立“并矢—指标—矩阵”三套记号的字典。核心只有两条作用规则；其余都是推论。读的时候请一直想着 §0 的主线：并矢是对象，矩阵是账本。

\subhdr{2.1\quad 阶：拿几个矢量喂它（必读）}

“张量”这个词让人紧张，但它的定义方式很朴素：\textbf{按它需要几个矢量才能给出一个数，来分阶。}

\begin{center}
\tabtop
\begin{tabular}{@{}p{0.12\linewidth}p{0.40\linewidth}p{0.40\linewidth}@{}}
\thead{阶} & \thead{吃什么、吐什么} & \thead{例子} \\
\multicolumn{3}{@{}l@{}}{\color{petal}\rule{\linewidth}{0.5pt}}\\
0 阶 & 不吃，直接是一个数 & 温度、质量、动能 \\
1 阶 & 吃一个矢量，吐一个数；且对输入线性 & 力在方向 $\vect u$ 上的投影 $\vect F\cdot\vect u$ \\
2 阶 & 吃两个矢量，吐一个数；且对每个输入都线性 & 应力：$\vect u\cdot\ten\sigma\cdot\vect v$；惯量：$\vect u\cdot\ten J\cdot\vect v$ \\
\end{tabular}
\tabbot
\end{center}

“对每个输入都线性”叫\textbf{多重线性}。对二阶就是要
\[
\ten T(\alpha\vect u_1+\beta\vect u_2,\;\vect v)=\alpha\ten T(\vect u_1,\vect v)+\beta\ten T(\vect u_2,\vect v),
\]
对第二个位置同样要求（两个位置各要求一次，合起来叫“双线性”）。

\textbf{为什么二阶张量有两个身份。}把双线性形式记作 $B(\vect u,\vect v)$。本稿固定\emph{第二个}位置：固定 $\vect v$ 后，$\vect u\mapsto B(\vect u,\vect v)$ 是“吃 $\vect u$ 吐数”的线性函数。在通常的三维空间里（有长度和点乘），每个这样的线性函数都能写成 $\vect w\cdot\vect u$，其中 $\vect w$ 唯一：取 $\vect w$ 为这个线性函数在三个基矢量上的值拼成的矢量即可。于是
\[
\vect v\;\longmapsto\;\bigl[\vect u\mapsto B(\vect u,\vect v)\bigr]
\]
是“吃 $\vect v$ 吐矢量 $\vect w$”的线性映射。反过来，给定线性映射 $\ten T$，令
\[
B(\vect u,\vect v)\stackrel{\text{def}}{=}(\ten T\vect v)\cdot\vect u ,
\]
它对两个位置都线性，所以是双线性形式。两者互逆，于是
\[
\boxed{\;\text{吃两个矢量吐一个数}\quad\Longleftrightarrow\quad\text{吃一个矢量吐一个矢量}\;}
\]
这两个说法是同一件事。课件 p9 写的是右边这个（$\vect b=\ten T\vect a$）；§1 用的是左边这个（$E=\vect a\cdot\ten T\vect a$）。\textbf{看起来是两套东西，其实是同一个对象的两个面。}

\textbf{槽位约定与转置（必读，后面全用这一条）。}上面特意固定\emph{第二}槽，是因为本稿的并矢右作用恰好这样做。若改固定\emph{第一}槽，得到的是同一台机器的转置。用一个非对称例子把两种识别都算一遍：取 $\vect a=\vect e_1$、$\vect b=\vect e_2$，
\[
B(\vect u,\vect v)=(\vect a\otimes\vect b)(\vect u,\vect v)=(\vect a\cdot\vect u)(\vect b\cdot\vect v).
\]
按本稿“固定第二槽”，对固定 $\vect v$，
\[
\vect u\mapsto(\vect a\cdot\vect u)(\vect b\cdot\vect v)=(\vect b\cdot\vect v)(\vect a\cdot\vect u)
=\bigl[\vect a(\vect b\cdot\vect v)\bigr]\cdot\vect u ,
\]
所以对应的机器是 $\vect v\mapsto\vect a(\vect b\cdot\vect v)=(\vect a\otimes\vect b)\cdot\vect v$。取输入 $\vect v=\vect e_1$，它给出 $\vect a(\vect b\cdot\vect e_1)=\vect0$。
若改固定第一槽，则对固定 $\vect u$，$\vect v\mapsto(\vect b\cdot\vect v)(\vect a\cdot\vect u)$ 对应的机器是 $\vect u\mapsto\vect b(\vect a\cdot\vect u)=(\vect b\otimes\vect a)\cdot\vect u$；同样输入 $\vect u=\vect e_1$，它给出 $\vect b$。两者相差一个转置（$\vect b\otimes\vect a=(\vect a\otimes\vect b)^{\mathsf T}$）。所以“输入 $\vect e_1$ 输出 $\vect e_2$ 还是 $\vect0$”取决于选了哪一槽；本稿之后一律固定第二槽，和 §2.2 的右作用规则一致，§2.3 的左作用则对应转置那一侧。

\subhdr{2.2\quad 张量积（并矢）：定义、作用、分量（必读）}

现在造一个具体的双线性对象。给两个矢量 $\vect a,\vect b$，规定
\[
(\vect a\otimes\vect b)(\vect u,\vect v)\stackrel{\text{def}}{=}(\vect a\cdot\vect u)(\vect b\cdot\vect v).
\]
右边是两个数的乘积，所以确实吐出一个数；对 $\vect u$ 线性、对 $\vect v$ 线性，都满足。这个对象记作 $\vect a\otimes\vect b$，叫 $\vect a$ 与 $\vect b$ 的\textbf{张量积}。\textbf{课件 p9 给它起的名字是 Dyad，中文叫并矢。}

它跟“喂一个矢量”的那个身份怎么对上？把第二个位置定死为 $\vect v$（这就是 §2.1 的槽位约定）：
\[
(\vect a\otimes\vect b)\cdot\vect v\;\stackrel{\text{def}}{=}\;\vect a\,(\vect b\cdot\vect v).
\]
这就是全部的作用规则，一句话：\textbf{先把 $\vect v$ 投影到 $\vect b$ 上，得到一个数，再拿这个数去放大 $\vect a$。}
另一边 $\vect v\cdot(\vect a\otimes\vect b)=(\vect v\cdot\vect a)\vect b$ 是另一条规则；它对应 §2.1 里“固定第一槽”的转置识别，二者一般不同，§2.3 专门讲。

\textbf{分量就是分量的乘积。}取正交基 $\{\vect e_i\}$，把 $i,j$ 各取 $1,2,3$，得到九个对象 $\vect e_i\otimes\vect e_j$。它们互相独立（不存在一组不全为零的系数使这些基并矢的组合为零），所以任何二阶张量都能唯一展开成
\[
\ten T=\sum_{i=1}^{3}\sum_{j=1}^{3}T_{ij}\,\vect e_i\otimes\vect e_j .
\]
“唯一”可以直接验证：把求和指标改名为 $i,l$，$\ten T=\sum_{i,l}T_{il}\,\vect e_i\otimes\vect e_l$，再作用在基矢量对 $(\vect e_j,\vect e_k)$ 上，只有 $i=j$、$l=k$ 那一项留下，得到 $T_{jk}$。所以九个系数就是九个分量，一个不多一个不少。
特别地
\[
(\vect a\otimes\vect b)_{ij}=a_ib_j
\qquad\Longrightarrow\qquad
\vect a\otimes\vect b=
\begin{pmatrix}
a_1b_1&a_1b_2&a_1b_3\\
a_2b_1&a_2b_2&a_2b_3\\
a_3b_1&a_3b_2&a_3b_3
\end{pmatrix}.
\]
矩阵的样子是“\textbf{列乘行}”，也叫外积。取 $\vect a=(1,2,0)$、$\vect b=(0,1,3)$，则
\[
\vect a\otimes\vect b=\begin{pmatrix}0&1&3\\0&2&6\\0&0&0\end{pmatrix}.
\]
\textbf{注意它不对称}：$(1,2)$ 元是 $1$，$(2,1)$ 元是 $0$。这就是“顺序有意义”的直接证据。再取 $\vect x=(1,1,1)$ 核对作用：
\[
(\vect a\otimes\vect b)\cdot\vect x=\vect a(\vect b\cdot\vect x)=\vect a\cdot4=(4,8,0),
\]
而矩阵乘 $\vect x$ 的三行分别是 $0+1+3=4$、$0+2+6=8$、$0$，一致。

\begin{center}
\begin{minipage}{\linewidth}\centering
\begin{tikzpicture}[scale=1.55,>={Stealth[length=1.7mm]}]
  \coordinate (O) at (0,0);
  \coordinate (A) at (1.0,1.2);
  \coordinate (B) at (1.0,0.0);
  \coordinate (X) at (0.4,1.0);
  \coordinate (R) at (0.4,0.48);
  \draw[inkgray,thin] (0,0)--(1.15,0);
  \draw[vv] (O)--(B) node[below]{\scriptsize $\vect b$};
  \draw[vv] (O)--(A) node[above left]{\scriptsize $\vect a$};
  \draw[vv] (O)--(X) node[left]{\scriptsize $\vect x$};
  \draw[inkgray,densely dashed] (X)--(0.4,0);
  \draw[inkgray] (0.4,0)--(0.46,0)--(0.46,0.06)--(0.4,0.06);
  \fill (0.4,0) circle (1pt);
  \node[inkgray,below] at (0.4,0) {\scriptsize $\vect b\cdot\vect x=0.4$};
  \draw[vviolet] (O)--(R) node[right]{\scriptsize $(\vect a\otimes\vect b)\cdot\vect x$};
  \draw[violet,densely dotted] (R)--(0.4,0);
\end{tikzpicture}\\[2pt]
\figcap{\mbox{\textbf{图 2}（自绘）：}并矢的右作用。先把 $\vect x$ 投影到 $\vect b$ 上，得到数 $\vect b\cdot\vect x=0.4$；再沿 $\vect a$ 放大同样的倍数，得到紫色箭头 $0.4\vect a$。图里 $\vect b$ 取单位矢量，所以投影长度就是点乘的值；一般情形只是多一个 $|\vect b|$ 的因子。}
\end{minipage}
\end{center}

\subhdr{2.3\quad 左右作用不一样（必读）}

并矢有两套作用规则，方向相反：
\[
\boxed{\;(\vect a\otimes\vect b)\cdot\vect x=\vect a\,(\vect b\cdot\vect x)\;}
\qquad\text{右作用：先投到 }\vect b\text{，再沿 }\vect a\text{ 放大};
\]
\[
\boxed{\;\vect x\cdot(\vect a\otimes\vect b)=(\vect x\cdot\vect a)\,\vect b\;}
\qquad\text{左作用：先投到 }\vect a\text{，再沿 }\vect b\text{ 放大}.
\]
\textbf{两者一般不同。}取 $\vect a=\vect e_1$、$\vect b=\vect e_2$、$\vect x=\vect e_1+2\vect e_2+3\vect e_3$：
\[
(\vect e_1\otimes\vect e_2)\cdot\vect x=\vect e_1(\vect e_2\cdot\vect x)=2\vect e_1,
\qquad
\vect x\cdot(\vect e_1\otimes\vect e_2)=(\vect x\cdot\vect e_1)\vect e_2=\vect e_2 .
\]
一个指向 $\vect e_1$，一个指向 $\vect e_2$，差别一眼可见。

\textbf{什么时候两者相同：当 $\ten D$ 对称时。}对一般二阶张量 $\ten D$，分量为
\[
(\ten D\cdot\vect x)_i=D_{ij}x_j,\qquad
(\vect x\cdot\ten D)_j=x_iD_{ij}=D_{ij}x_i=(D^{\mathsf T}\vect x)_j .
\]
中间那个等号只是把两个数相乘的顺序换了一下；最后一步用了 $(D^{\mathsf T})_{ji}=D_{ij}$，即 $(\ten D^{\mathsf T}\vect x)_j=(D^{\mathsf T})_{ji}x_i=D_{ij}x_i$。所以 $\ten D\cdot\vect x=\vect x\cdot\ten D$ 对\textbf{所有} $\vect x$ 成立，当且仅当 $\ten D=\ten D^{\mathsf T}$（逐分量就是 $D_{ij}=D_{ji}$）。力学里的 $\ten J$、$\ten\sigma$ 都对称，所以“碰巧”怎么写都对；但 §3.2 的旋转算子 $S$ \textbf{一般不对称}（§3.2 会给出它的定义），左右一混立刻出错。

\begin{center}
\fcolorbox{violet}{bandviolet}{\parbox{0.94\linewidth}{\centering\vspace{3pt}
\small 一句话：$\vect a\otimes\vect b$ 一般不对称，$\vect a\otimes\vect b+\vect b\otimes\vect a$ 一定对称；当 $\vect a$ 与 $\vect b$ 平行时（例如 $\vect a=\vect b=\vect n$），$\vect n\otimes\vect n$ 本身就是对称的（§2.5 的投影并矢）。并矢带着顺序信息，所以它不是“两个矢量的乘积”。\vspace{3pt}}}
\end{center}

\subhdr{2.4\quad 点乘、双点积、迹（必读）}

两个并矢之间可以“乘”，乘掉几个指标就叫几点积：
\[
(\vect a\otimes\vect b)\cdot(\vect c\otimes\vect d)=(\vect b\cdot\vect c)\,\vect a\otimes\vect d
\qquad\text{（缩并掉一个指标，剩二阶）},
\]
\[
(\vect a\otimes\vect b):(\vect c\otimes\vect d)=(\vect a\cdot\vect c)(\vect b\cdot\vect d)
\qquad\text{（缩并掉两个指标，剩标量）}.
\]
第二条叫\textbf{双点积}，写成指标就是 $\ten A:\ten B=\sum_{i,j}A_{ij}B_{ij}$，与矩阵写法 $\tr(\ten A^{\mathsf T}\ten B)$ 等价：因为 $(\ten A^{\mathsf T}\ten B)_{jj}=A_{ij}B_{ij}$，再对 $j$ 求和。

\textbf{迹。}$\tr\ten D=\sum_iD_{ii}$。用基并矢写，$\tr(\vect a\otimes\vect b)=\vect a\cdot\vect b$；于是
\[
\unitd:\ten D=\tr\ten D,\qquad
\ten A:(\vect u\otimes\vect v)=\vect u\cdot\ten A\cdot\vect v .
\]
\textbf{数字例子。}取 $\ten D=\vect e_1\otimes\vect e_1+2\vect e_1\otimes\vect e_2+3\vect e_2\otimes\vect e_2$，矩阵是 $\begin{pmatrix}1&2\\0&3\end{pmatrix}$（二维示意，三维同理）。则
\[
\tr\ten D=1+3=4,\qquad
\ten D:\ten D=1^2+2^2+3^2=14,\qquad
\unitd:\ten D=4 .
\]
最后一条可以当公式用：$\unitd$ 与谁双点积就取谁的迹。

\subhdr{2.5\quad 单位并矢与投影并矢（必读）}

\textbf{单位并矢}就是克罗内克符号那个张量：
\[
\unitd=\sum_{k=1}^{3}\vect e_k\otimes\vect e_k
\quad\Longrightarrow\quad
\unitd\cdot\vect x=\sum_k\vect e_k(\vect e_k\cdot\vect x)=\vect x
\qquad\text{（对任意正交基成立）}.
\]
最后一个等号就是“把 $\vect x$ 在三个正交方向上的投影拼回来”。$\unitd$ 换基后还是 $\unitd$，§2.7 会用变换律再证一次。

\textbf{投影并矢。}设 $\vect n$ 是单位矢量，令 $\ten P=\vect n\otimes\vect n$，则
\[
\ten P\cdot\vect x=\vect n(\vect n\cdot\vect x)
\quad\text{（$\vect x$ 在 }\vect n\text{ 上的投影）},
\qquad
\ten P^2=\ten P\ \text{（再投一次不变）}.
\]
而 $\unitd-\vect n\otimes\vect n$ 是“投影到垂直于 $\vect n$ 的平面”的算子。两者的迹分别是 $1$ 与 $2$：$\ten P$ 的迹为 $1$，对应它\textbf{保留}沿 $\vect n$ 的 $1$ 个方向、去掉垂直平面里的 $2$ 个；$\unitd-\ten P$ 的迹为 $2$，对应它\textbf{保留}垂直平面里的 $2$ 个方向、去掉沿 $\vect n$ 的 $1$ 个。在三维里，一个投影算子的迹正好等于它保留的方向数（把“投影掉几个自由度”说反，就会把保留维数和消去维数颠倒）。

\begin{center}
\begin{tikzpicture}[scale=1.8,>={Stealth[length=1.7mm]}]
  \coordinate (O) at (0,0);
  \coordinate (N) at (0.866,0.5);
  \coordinate (X) at (0.2,1.1);
  \coordinate (P) at (0.6263,0.3616);
  \draw[inkgray,densely dashed] (-0.1,-0.058)--(1.05,0.606);
  \draw[vv] (O)--(N) node[above right]{\scriptsize $\vect n$};
  \draw[vv] (O)--(X) node[left]{\scriptsize $\vect x$};
  \draw[vviolet] (O)--(P) node[below right]{\scriptsize $\ten P\cdot\vect x$};
  \draw[inkgray,thick,->,>={Stealth[length=1.7mm]}] (P)--(X) node[right]{\scriptsize $(\unitd-\ten P)\cdot\vect x$};
  \draw[inkgray,thin] (0.6783,0.3916)--(0.6483,0.4436)--(0.5963,0.4136);
  \fill (O) circle (0.7pt);
\end{tikzpicture}\\[2pt]
\figcap{\mbox{\textbf{图 3}（自绘）：}投影并矢 $\ten P=\vect n\otimes\vect n$ 把 $\vect x$ 压到 $\vect n$ 所在的直线上，得到紫箭头 $\ten P\cdot\vect x=\vect n(\vect n\cdot\vect x)$；剩下的灰箭头是 $(\unitd-\ten P)\cdot\vect x=\vect x-\vect n(\vect n\cdot\vect x)$，它落在与 $\vect n$ 垂直的平面里（图中直角标记）。两个分量之和就是 $\vect x$ 本身——这就是 $\ten P+(\unitd-\ten P)=\unitd$ 的几何意思。\textbf{这里的 $\unitd$ 是单位并矢（二阶张量），不是数 $1$}。}
\end{center}

\subhdr{2.6\quad 转置、对称分解、对偶张量与叉乘（必读）}

\textbf{转置。}
\[
(\vect a\otimes\vect b)^{\mathsf T}=\vect b\otimes\vect a,
\qquad
(\ten A\cdot\ten B)^{\mathsf T}=\ten B^{\mathsf T}\cdot\ten A^{\mathsf T},
\qquad
(\ten D^{\mathsf T})_{ij}=D_{ji}.
\]
转置的判别式很好用：\textbf{对任意 $\vect x,\vect y$ 有 $(\ten D\cdot\vect x)\cdot\vect y=\vect x\cdot(\ten D^{\mathsf T}\cdot\vect y)$}。§1.5 证特征向量正交时用的就是这一条。

\textbf{对称分解。}任何并矢都能唯一拆成对称与反对称两半：
\[
\ten D=\underbrace{\tfrac12(\ten D+\ten D^{\mathsf T})}_{\text{对称}}\;+\;
\underbrace{\tfrac12(\ten D-\ten D^{\mathsf T})}_{\text{反对称}} .
\]
力学里这两半各有明确身份：\textbf{对称半}给惯量、应力、应变；\textbf{反对称半}等价于一个矢量，正是叉乘（下面马上说）。

\textbf{叉乘的紧凑写法。}定义\textbf{列维—奇维塔符号} $\varepsilon_{ijk}$：若 $(i,j,k)$ 是 $(1,2,3)$ 的偶排列（$123,231,312$）取 $+1$；奇排列（$132,213,321$）取 $-1$；有重复指标取 $0$。“偶/奇排列”就是把 $(1,2,3)$ 经过偶数次/奇数次两两交换得到；六个排列已在括号里列全，直接查表即可。于是
\[
(\vect a\times\vect b)_i=\varepsilon_{ijk}a_jb_k .
\]
核对第一分量：$\varepsilon_{123}a_2b_3+\varepsilon_{132}a_3b_2=a_2b_3-a_3b_2$，与 §0.1 的分量式一致。

\textbf{对偶张量。}对任意矢量 $\vect a$，定义二阶张量 $\vect a^{\times}$ 的分量为
\[
(\vect a^{\times})_{ij}=\varepsilon_{ikj}a_k,
\qquad
\vect a^{\times}=
\begin{pmatrix}
0&-a_3&a_2\\
a_3&0&-a_1\\
-a_2&a_1&0
\end{pmatrix}.
\]
它的作用正好是叉乘：
\[
(\vect a^{\times}\cdot\vect x)_i=\varepsilon_{ikj}a_kx_j=\varepsilon_{ijk}a_jx_k=(\vect a\times\vect x)_i,
\]
即 $\vect a^{\times}\cdot\vect x=\vect a\times\vect x$。反过来，若 $\ten A$ 反对称，则存在唯一的矢量 $\vect a$ 使 $\ten A=\vect a^{\times}$，反解公式是
\[
a_k=-\tfrac12\varepsilon_{kij}A_{ij}.
\]
这就是课件 p8 的对偶张量。\textbf{一句话：叉乘不是“另一种乘法”，而是一个反对称并矢在作用。}
“为什么叉乘不满足结合律”也就通了：不满足结合律的是叉乘这个运算；而反对称张量的\textbf{合成}（矩阵乘法）是满足结合律的。

\textbf{旋转算子的并矢写法。}有了对偶张量，§2.5 的两个投影并矢就能拼出课件 p11 的旋转算子：
\[
S(\varphi;\vect n)=\vect n\otimes\vect n+(\unitd-\vect n\otimes\vect n)\cos\varphi+\vect n^{\times}\sin\varphi .
\]
读法：\textbf{沿轴的分量不动（第一项），垂直轴的分量被转了 $\cos$ 又在 $\vect n\times$ 方向上补了 $\sin$（后两项）。}这个公式 §3.4 会逐项验证；现在只需看出：它用了单位并矢、投影并矢和对偶张量三样东西。

\warnbox{\textbf{符号约定要看清楚。}课件 p8 写的是 $a^{\times}_{ik}=\varepsilon_{ijk}a_j$，与本稿 $(\vect a^{\times})_{ij}=\varepsilon_{ikj}a_k$ 是同一个矩阵（把哑指标改个名字即可）；RIT Whelan notes 10 的 p7 用的是 $\ast A\cdot\vect B=-\vect A\times\vect B$，即它的 $\ast A$ 等于 $-\vect A^{\times}$。读不同书时先认矩阵的三个元素，别只记一个公式。}

\subhdr{2.7\quad 商法则：怎样确认一组九个数是张量（必读）}

随便给我九个排成方阵的数，它是不是张量？\textbf{不一定。}张量的定义不在“有九个分量”，而在“换基时这九个分量怎么变”。规则就是 §0.4 那条：
\[
\boxed{\;T'_{ij}=R_{ki}R_{lj}T_{kl}\;}
\qquad\text{（$R_{ki}=\vect n_i\cdot\vect e_k$ 是新基在旧基下的分量）}
\]
凡按这条规则变的九元数组，才是张量；不按这条变的，只是九个数字。

实践中更常用的是反过来：手上有一个物理关系，想证明里面那个九元数组是张量。这就是\textbf{商法则}：

\begin{center}
\fcolorbox{violet}{bandviolet}{\parbox{0.94\linewidth}{\vspace{3pt}
\textbf{商法则（三维，实用版）。}若对\textbf{任意}矢量 $\vect v$ 都有 $\vect u=\ten T\vect v$（即 $u_i=T_{ij}v_j$），且已知 $\vect u$、$\vect v$ 都是矢量（换基时按 $\vect u'=\ten R^{\mathsf T}\vect u$、$\vect v'=\ten R^{\mathsf T}\vect v$ 变），则 $\ten T$ 必是二阶张量，即 $\ten T'=\ten R^{\mathsf T}\ten T\ten R$。
\vspace{3pt}}}
\end{center}

证明只有三步。在带撇的坐标系里同一条定律必须同样成立，$\vect u'=\ten T'\vect v'$；把 $\vect u'=\ten R^{\mathsf T}\vect u$、$\vect v'=\ten R^{\mathsf T}\vect v$ 代进去：
\[
\ten R^{\mathsf T}\vect u=\ten T'\ten R^{\mathsf T}\vect v
\quad\Longrightarrow\quad
\vect u=\ten R\ten T'\ten R^{\mathsf T}\vect v .
\]
与原来的 $\vect u=\ten T\vect v$ 对任意 $\vect v$ 比较，得 $\ten T=\ten R\ten T'\ten R^{\mathsf T}$，两边左乘 $\ten R^{\mathsf T}$、右乘 $\ten R$（用 $\ten R^{\mathsf T}\ten R=\unitd$）就得到 $\ten T'=\ten R^{\mathsf T}\ten T\ten R$。
\textbf{所以“它是张量”不是假设，是被这条定律逼出来的。}§6.1 会用它：面力 $\vect t$ 与法向 $\vect n$ 都是矢量，而 $\vect t(\vect n)=\ten\sigma\cdot\vect n$ 对任意 $\vect n$ 成立，于是 $\ten\sigma$ 一定是二阶张量。

\textbf{一个真的张量：$\delta_{ij}$。}把 $T_{kl}=\delta_{kl}$ 代进变换律：
\[
\delta'_{ij}=R_{ki}R_{lj}\delta_{kl}=R_{ki}R_{kj}=(\ten R^{\mathsf T}\ten R)_{ij}=\delta_{ij}.
\]
所以单位张量在任何正交基下的分量都是 $\delta_{ij}$——这正是“单位矩阵换基后还是单位矩阵”的深层原因，也是 §2.5 的 $\unitd=\sum_k\vect e_k\otimes\vect e_k$ 对任意正交基都成立的理由。

\textbf{一个不是张量的东西：位置矢量。}三个数 $(x_1,x_2,x_3)$ 在\textbf{旋转}下确实按矢量规则变（所以转动问题里可以当矢量用），但一旦允许\textbf{平移}坐标系，它会多出一个常数项，不再是齐次的（“齐次”就是分量只含坐标的一次项、不含常数项）。所以“是不是张量”这句话必须连着\textbf{允许哪些变换}说。本稿（以及整门理论力学）只允许旋转与反射，即\textbf{正交变换}，在这个群里位置矢量是矢量。

\subhdr{2.8\quad 缩并层级表（查阅）}

这张表只是把 §2.2--§2.7 的规则按“缩并掉几个指标”读了一遍。\textbf{它不是要你背，而是做题时横向对照。}

\begin{center}
\tabtop
\begin{tabular}{@{}lll@{}}
\thead{写法} & \thead{结果} & \thead{分量} \\
\multicolumn{3}{@{}l@{}}{\color{petal}\rule{\linewidth}{0.5pt}}\\
$\ten A\cdot\ten B$ & 二阶张量 & $A_{ik}B_{kj}$ \\
$\ten A\cdot\vect x$ & 矢量 & $A_{ij}x_j$ \\
$\vect x\cdot\ten A$ & 矢量 & $x_iA_{ij}$ \\
$\ten A:\ten B$ & 标量 & $A_{ij}B_{ij}$ \\
$\tr\ten A$ & 标量 & $A_{ii}$ \\
$(\vect a\otimes\vect b)\cdot(\vect c\otimes\vect d)$ & 二阶张量 & $(\vect b\cdot\vect c)\,\vect a\otimes\vect d$ \\
$\vect a\cdot(\vect b\otimes\vect c)$ & 矢量 & $(\vect a\cdot\vect b)\vect c$ \\
$(\vect a\otimes\vect b):(\vect c\otimes\vect d)$ & 标量 & $(\vect a\cdot\vect c)(\vect b\cdot\vect d)$ \\
\end{tabular}
\tabbot
\end{center}

逐行读法：前五行的规律是“每点一下吃掉一个指标”，所以二阶$\cdot$二阶留二阶、二阶$\cdot$矢量留一阶、双点积吃掉两个留零阶。
最后三行要小心，它们是整表最容易错的地方：
\begin{itemize}
\item 第六行：两个并矢点乘，缩并的是\textbf{右并矢的左因子}与\textbf{左并矢的右因子}，即 $\vect b$ 与 $\vect c$。这一条是 §3.3 证明 $S^{\mathsf T}S=\unitd$ 的关键一步。
\item 第七行：矢量点并矢，$\vect a\cdot(\vect b\otimes\vect c)=(\vect a\cdot\vect b)\vect c$，结果是矢量，方向沿 $\vect c$。它和第六行差在“谁先被缩并”，不要凭形状猜。
\item 第八行：双点积把两个并矢的四个因子两两配对，结果是一个数；它等于矩阵写法 $\tr(\ten A^{\mathsf T}\ten B)$。
\end{itemize}

\subhdr{2.9\quad 谱分解：把并矢搬到主轴上（必读）}

§1.5 的谱定理用并矢写出来，就是一句话：
\[
\ten D=\lambda_1\,\vect n_1\otimes\vect n_1+\lambda_2\,\vect n_2\otimes\vect n_2+\lambda_3\,\vect n_3\otimes\vect n_3 .
\]
这三项每一项都是一个\textbf{投影并矢}乘上系数，所以“沿主轴分解”的物理含义是：\textbf{把任意矢量先投影到三个主轴上，各自乘 $\lambda_k$，再拼回来。}

\textbf{用 §1.4 的例子核对。}那里 $\ten D=\begin{pmatrix}6&2&0\\2&6&0\\0&0&1\end{pmatrix}$，$\lambda_1=8,\lambda_2=4,\lambda_3=1$，$\vect n_1=(1,1,0)/\sqrt2$、$\vect n_2=(1,-1,0)/\sqrt2$、$\vect n_3=\vect e_3$。取 $\vect x=(1,2,3)$，三个投影是
\[
\vect n_1\cdot\vect x=\frac{3}{\sqrt2},\qquad
\vect n_2\cdot\vect x=-\frac{1}{\sqrt2},\qquad
\vect n_3\cdot\vect x=3 .
\]
于是
\[
\ten D\vect x=8\cdot\frac{3}{\sqrt2}\vect n_1+4\cdot\Bigl(-\frac{1}{\sqrt2}\Bigr)\vect n_2+1\cdot3\,\vect e_3
=(12,12,0)+(-2,2,0)+(0,0,3)=(10,14,3),
\]
与直接做矩阵乘法（第一行 $6+4=10$，第二行 $2+12=14$，第三行 $3$）一致。

\textbf{椭球图像（接 §1.3）。}对对称正定的 $\ten D$，$\vect x\cdot\ten D\cdot\vect x=1$ 是椭球，三个半轴 $1/\sqrt{\lambda_k}$，方向就是主轴。惯量椭球在惯量张量正定时是这个图像；应力张量的主值若不同号，等值面不是椭球，这时只用主轴与主应力，不画椭球。若某个 $\lambda=0$，等值面退化成柱面或平面——这意味着那台机器在某个方向上“什么都不给”，$\det\ten D=0$（§0.5）。

\textbf{三个不变量（接 §1.6）}在并矢语言里同样成立：$I_1=\tr\ten D$、$I_3=\det\ten D$、$I_2=\tfrac12[(\tr\ten D)^2-\tr\ten D^2]$，做题时用来验算。

\subhdr{2.10\quad 三记号字典：并矢、指标、矩阵（查阅）}

下表把同一个对象在三种写法下并排列出。\textbf{它不是要你背，而是你读到任何一本书时可以横向对照。}

\begin{center}
\begin{minipage}{\linewidth}
\small
\renewcommand{\arraystretch}{1.12}
\tabtop
\begin{tabular}{@{}p{0.28\linewidth}p{0.20\linewidth}p{0.36\linewidth}@{}}
\thead{Gibbs 并矢} & \thead{指标} & \thead{矩阵} \\
\multicolumn{3}{@{}l@{}}{\color{petal}\rule{\linewidth}{0.5pt}}\\
$\vect a\vect b$ 或 $\vect a\otimes\vect b$ & $a_ib_j$ & $\vect a\,\vect b^{\mathsf T}$（列乘行） \\
$\ten D\cdot\vect x$ & $D_{ij}x_j$ & $\ten D\,\vect x$ \\
$\vect x\cdot\ten D$ & $x_iD_{ij}$ & $\vect x^{\mathsf T}\ten D$ \\
$\ten A\cdot\ten B$ & $A_{ik}B_{kj}$ & $\ten A\ten B$ \\
$\ten A:\ten B$ & $A_{ij}B_{ij}$ & $\tr(\ten A^{\mathsf T}\ten B)$ \\
$\ten D^{\mathsf T}$ & $D_{ji}$ & $\ten D^{\mathsf T}$ \\
$\unitd$ & $\delta_{ij}$ & 单位矩阵 \\
$\tr\ten D$ & $D_{ii}$ & 对角元之和 \\
$\vect a^{\times}$ & $\varepsilon_{ikj}a_k$ & 反对称矩阵（§2.6） \\
$\ten D=\sum_k\lambda_k\vect n_k\vect n_k^{\mathsf T}$ & $\sum_k\lambda_k n_{ki}n_{kj}$ & $\ten R\operatorname{diag}(\lambda)\ten R^{\mathsf T}$ \\
\end{tabular}
\tabbot
\end{minipage}
\end{center}

\textbf{表的最后一行值得单独看。}它就是“对角化”的全部内容：找一组正交基，使矩阵变成对角阵。
矩阵写法要写三个东西（$\ten R$、对角线、$\ten R^{\mathsf T}$），并矢写法只写一项求和，而物理含义在并矢写法里是\textbf{看得见}的：三个投影并矢各自乘上主值。这就是为什么力学书偏爱并矢。

\textbf{但并矢不是矩阵。}矩阵是“坐标下的账本”，并矢是不依赖坐标系的对象。换基时账本变、对象不变（§0.4 的主线）。你可以在任意基下把并矢写成九个分量，但并矢本身不是那九个数——§2.7 已经说明，只有按变换律变的九个数才配叫张量。

\warnbox{\textbf{选读（可跳过）：上下标是从哪来的。}
你以后会看到 $T^{i}{}_{j}$、$T_{ij}$、$T^{ij}$ 三种写法并存。在\textbf{正交}基里，三种写法数值完全相同，所以随便挑一种。但在\textbf{非正交}基（斜坐标、曲线坐标）里，矢量有两套分量：把矢量投影到基矢量上得到一套（协变，下标记），把矢量沿基矢量方向分解的系数得到另一套（逆变，上标记）；二者由一个叫度规的 $g_{ij}$ 互相转换，$g_{ij}$ 就是基矢量的点乘表 $g_{ij}=\vect e_i\cdot\vect e_j$。只有当 $g_{ij}=\delta_{ij}$（正交归一）时，两套才重合，上下标才可以不分。
\textbf{结论：}理论力学的刚体部分全程正交基，所以本稿一律用下标，不会错。等到连续介质力学用曲线坐标、或相对论用四维时空时，上下标就必须写清楚——那时再回来补这半页即可。跳过后从 §3 继续。}

\hdr{§3}{旋转算子与角速度：课件 p10--p12 的并矢读法}

这一节把课件的旋转算子 $S$ 讲成“把旧基送到新基的那台机器”，再从它的导数得到角速度。它同时是并矢运算的一次综合练习：张量积、单位并矢、投影并矢、对偶张量都会出现。

\subhdr{3.1\quad 两个基与一个算子（必读）}

取固定系的正交基 $\vect e_1,\vect e_2,\vect e_3$（课件 p10 记作 $\vect n_1,\vect n_2,\vect n_3$），再取随刚体转动的正交基 $\vect n'_1(t),\vect n'_2(t),\vect n'_3(t)$，并约定
\[
\vect n'_k(0)=\vect e_k,\qquad k=1,2,3 .
\]
刚体上任意一点的位置矢量可以写成
\[
\vect r(t)=x'_1\vect n'_1(t)+x'_2\vect n'_2(t)+x'_3\vect n'_3(t)=x'_k\vect n'_k(t),
\]
其中 $x'_k$ 是\textbf{随体坐标}，对刚体上的点不随时间变（刚体上任意两点的距离不变，所以随体坐标不变）。在 $t=0$ 时 $\vect r(0)=x'_k\vect e_k$。

我们要找一台“机器” $S(t)$，它把 $t=0$ 的位置送到 $t$ 时刻的位置：
\[
\vect r(t)=S(t)\vect r(0),\qquad\text{并且}\quad \vect n'_k(t)=S(t)\vect e_k .
\]
第二句是第一句的“基矢量版本”：如果 $S$ 把每个旧基矢量都送到对应的新基矢量，那么由线性，它就把任意矢量送到正确的位置。

\subhdr{3.2\quad $S(t)=\sum_k\vect n'_k\otimes\vect e_k$（必读）}

课件 p10 直接给出
\[
S(t)=\sum_{k=1}^{3}\vect n'_k(t)\otimes\vect e_k .
\]
为什么是它？把它作用在旧基矢量 $\vect e_j$ 上：
\[
S(t)\vect e_j=\sum_k\vect n'_k(\vect e_k\cdot\vect e_j)=\vect n'_j .
\]
正好把 $\vect e_j$ 送到 $\vect n'_j$。再作用在 $\vect r(0)=x'_k\vect e_k$ 上：
\[
S(t)\vect r(0)=x'_kS(t)\vect e_k=x'_k\vect n'_k(t)=\vect r(t).
\]
所以 $S$ 的定义不是猜的，是“要求 $S\vect e_k=\vect n'_k$”唯一确定下来的；写成并矢和只是把这个要求一次性满足。

\textbf{矩阵读法。}$S$ 在固定基下的矩阵，第 $k$ 列就是 $\vect n'_k$ 的三个分量（因为 $S\vect e_k=\vect n'_k$）。二维绕 $\vect e_3$ 转 $\varphi$ 的例子：
\[
\vect n'_1=\cos\varphi\,\vect e_1+\sin\varphi\,\vect e_2,\qquad
\vect n'_2=-\sin\varphi\,\vect e_1+\cos\varphi\,\vect e_2,
\]
\[
S(\varphi;\vect e_3)=
\begin{pmatrix}
\cos\varphi&-\sin\varphi\\
\sin\varphi&\cos\varphi
\end{pmatrix},
\qquad
S\vect e_1=\vect n'_1,\quad S\vect e_2=\vect n'_2 .
\]
这正是 §0.4 里那个 $R$（列是新基分量的矩阵）。所以 §0.4 的换基矩阵和这里的旋转算子是同一个东西；只是那里用它换账本，这里用它转矢量。

\begin{center}
\begin{tikzpicture}[scale=1.75,>={Stealth[length=1.7mm]}]
  \coordinate (O) at (0,0);
  \coordinate (E1) at (1,0);
  \coordinate (E2) at (0,1);
  \coordinate (N1) at (0.8192,0.5736);
  \coordinate (N2) at (-0.5736,0.8192);
  \draw[ax] (O)--(1.25,0) node[right]{\scriptsize $\vect e_1$};
  \draw[ax] (O)--(0,1.25) node[above]{\scriptsize $\vect e_2$};
  \draw[vviolet] (O)--(N1) node[right]{\scriptsize $\vect n'_1$};
  \draw[vviolet] (O)--(N2) node[above left]{\scriptsize $\vect n'_2$};
  \draw[inkgray,thin] (0.48,0) arc (0:35:0.48);
  \node[inkgray] at (0.60,0.16) {\scriptsize $\varphi$};
  \draw[inkgray,thin] (0,0) circle (0.075);
  \fill (0,0) circle (0.025);
  \node[inkgray,below right] at (0.04,-0.02) {\scriptsize $\vect n$（出纸面）};
\end{tikzpicture}\\[2pt]
\figcap{\mbox{\textbf{图 4}（据课件 p10 的定义自绘）：}在垂直于转轴 $\vect n$ 的平面里，旧基 $\vect e_1,\vect e_2$（灰虚线）被 $S$ 送到新基 $\vect n'_1,\vect n'_2$（紫线），转角 $\varphi$。转轴 $\vect n$ 指向纸外（$\odot$）。$S$ 的矩阵第 $1$ 列是 $\vect n'_1$ 的分量，第 $2$ 列是 $\vect n'_2$ 的分量；$S\vect e_k=\vect n'_k$ 是这张图的公式版本。}
\end{center}

\subhdr{3.3\quad $S$ 是正交算子：长度与点乘都不变（必读）}

用 §2.8 的并矢点乘规则算 $S^{\mathsf T}S$。先写
\[
S=\sum_k\vect n'_k\otimes\vect e_k,\qquad
S^{\mathsf T}=\sum_l\vect e_l\otimes\vect n'_l .
\]
于是
\[
S^{\mathsf T}S=\sum_{k,l}(\vect e_l\otimes\vect n'_l)\cdot(\vect n'_k\otimes\vect e_k)
=\sum_{k,l}(\vect n'_l\cdot\vect n'_k)\,\vect e_l\otimes\vect e_k
=\sum_k\vect e_k\otimes\vect e_k=\unitd .
\]
中间那步用了 $\vect n'_l\cdot\vect n'_k=\delta_{lk}$（新基也是正交归一的）。同法
\[
SS^{\mathsf T}=\sum_{k,l}(\vect n'_k\otimes\vect e_k)\cdot(\vect e_l\otimes\vect n'_l)
=\sum_{k,l}(\vect e_k\cdot\vect e_l)\,\vect n'_k\otimes\vect n'_l
=\sum_k\vect n'_k\otimes\vect n'_k=\unitd .
\]
所以 $S^{-1}=S^{\mathsf T}$。它保持点乘：
\[
(S\vect x)\cdot(S\vect y)=\vect x\cdot(S^{\mathsf T}S\vect y)=\vect x\cdot\vect y .
\]
长度、夹角都不变，这正是“刚体转动”该有的性质。又因为 $\det(S^{\mathsf T}S)=(\det S)^2=1$，所以 $\det S=\pm1$；$S(t)$ 随 $t$ 连续变化且 $S(0)=\unitd$ 的 $\det$ 为 $+1$，所以\textbf{连续转动给 $\det S=+1$}（$\det=-1$ 对应镜像反射，比如把右手变成左手，不是刚体转动）。

\subhdr{3.4\quad 绕轴转 $\varphi$ 的显式公式（必读）}

课件 p11 给出
\[
S(\varphi;\vect n)=\vect n\otimes\vect n+(\unitd-\vect n\otimes\vect n)\cos\varphi+\vect n^{\times}\sin\varphi .
\]
逐项验证。把任意矢量 $\vect x$ 拆成“沿轴”与“垂直轴”两部分：
\[
\vect x=\vect x_\parallel+\vect x_\perp,\qquad
\vect x_\parallel=(\vect n\cdot\vect x)\vect n=(\vect n\otimes\vect n)\vect x,\qquad
\vect x_\perp=\vect x-\vect x_\parallel .
\]
于是
\[
(\vect n\otimes\vect n)\vect x=\vect x_\parallel,\qquad
(\unitd-\vect n\otimes\vect n)\vect x=\vect x_\perp,\qquad
\vect n^{\times}\vect x=\vect n\times\vect x=\vect n\times\vect x_\perp
\]
（最后一步因为 $\vect n\times\vect x_\parallel=0$）。代入公式：
\[
S\vect x=\vect x_\parallel+\vect x_\perp\cos\varphi+(\vect n\times\vect x_\perp)\sin\varphi .
\]
读法：\textbf{沿轴分量不动；垂直轴分量在垂直于 $\vect n$ 的平面里转了 $\varphi$（$\cos$ 项是原方向，$\sin$ 项是转 $90^\circ$ 后的方向）。}这就是右手定则的转动。

\textbf{两个核对。}取 $\vect n=\vect e_3$、$\varphi=90^\circ$、$\vect x=\vect e_1$：
\[
S\vect e_1=\vect e_3\otimes\vect e_3\vect e_1+0\cdot(\unitd-\vect e_3\otimes\vect e_3)\vect e_1+\vect e_3^{\times}\vect e_1=\vect e_2,
\]
与“绕 $z$ 轴把 $x$ 转到 $y$”一致。取 $\varphi=0$：$S=\vect n\otimes\vect n+(\unitd-\vect n\otimes\vect n)=\unitd$，也对。

\textbf{注意 $S$ 一般不对称。}$\vect n\otimes\vect n$ 与 $\unitd-\vect n\otimes\vect n$ 都对称，但 $\vect n^{\times}$ 反对称，所以 $S$ 的转置是
\[
S^{\mathsf T}(\varphi;\vect n)=\vect n\otimes\vect n+(\unitd-\vect n\otimes\vect n)\cos\varphi-\vect n^{\times}\sin\varphi=S(-\varphi;\vect n).
\]
因此 $S\cdot\vect x$ 与 $\vect x\cdot S$ 一般不同（§2.3 的警告在这里兑现）：前者是“把矢量转 $+\varphi$”，后者是“把矢量转 $-\varphi$”。

\subhdr{3.5\quad 无穷小旋转与角速度（必读）}

把 $\varphi$ 取成无穷小 $\dd\varphi$，用 $\cos\dd\varphi\approx1$、$\sin\dd\varphi\approx\dd\varphi$（只保留 $\dd\varphi$ 的一次项，忽略 $\dd\varphi^2$ 及更高次）：
\[
S(\dd\varphi;\vect n)=\vect n\otimes\vect n+(\unitd-\vect n\otimes\vect n)+\vect n^{\times}\dd\varphi
=\unitd+\dd\varphi\,\vect n^{\times}=\unitd+\dd{\boldsymbol\varphi}^{\times},
\]
其中 $\dd{\boldsymbol\varphi}=\dd\varphi\,\vect n$ 是“转轴乘转角”的矢量。两个无穷小旋转相乘：
\[
S_1S_2=(\unitd+\dd{\boldsymbol\varphi}_1^{\times})(\unitd+\dd{\boldsymbol\varphi}_2^{\times})
=\unitd+\dd{\boldsymbol\varphi}_1^{\times}+\dd{\boldsymbol\varphi}_2^{\times}+O(\dd\varphi^2),
\]
一阶项对易（相乘顺序可以交换），所以\textbf{无穷小转动可交换}；有限转动一般不可交换（本稿不用这个结论）。

现在对 $S(t)S^{\mathsf T}(t)=\unitd$ 求导：
\[
\frac{\dd S}{\dd t}S^{\mathsf T}+S\frac{\dd S^{\mathsf T}}{\dd t}=0 .
\]
记
\[
\ten A=\frac{\dd S}{\dd t}S^{\mathsf T}.
\]
则
\[
\ten A^{\mathsf T}=\Bigl(\frac{\dd S}{\dd t}S^{\mathsf T}\Bigr)^{\mathsf T}
=S\frac{\dd S^{\mathsf T}}{\dd t}=-\ten A,
\]
所以 $\ten A$ 反对称。由 §2.6，反对称二阶张量恰好对应一个矢量，记作 $\boldsymbol\Omega$，使
\[
\boxed{\;\boldsymbol\Omega^{\times}=\frac{\dd S}{\dd t}S^{\mathsf T}\;}
\]
这就是课件 p12 的角速度定义。$\boldsymbol\Omega$ 叫\textbf{角速度矢量}；它不是“九个分量”，而是反对称张量的三个独立分量（§2.6），所以是一个矢量。

\textbf{为什么 $\dd\vect n'_k/\dd t=\boldsymbol\Omega\times\vect n'_k$。}由 $\ten A=\boldsymbol\Omega^{\times}$ 与 $S^{\mathsf T}S=\unitd$，有
\[
\frac{\dd S}{\dd t}=\boldsymbol\Omega^{\times}S .
\]
两边作用在 $\vect e_k$ 上，并用 $S\vect e_k=\vect n'_k$：
\[
\frac{\dd\vect n'_k}{\dd t}=\frac{\dd}{\dd t}(S\vect e_k)=\frac{\dd S}{\dd t}\vect e_k=\boldsymbol\Omega^{\times}S\vect e_k=\boldsymbol\Omega\times\vect n'_k .
\]
这正是“随体基矢量以角速度 $\boldsymbol\Omega$ 转动”的微分版本。作为核对，也可以直接用并矢算：
\[
\frac{\dd S}{\dd t}S^{\mathsf T}
=\sum_k\frac{\dd\vect n'_k}{\dd t}\otimes\vect n'_k,
\qquad
\boldsymbol\Omega^{\times}=\sum_k(\boldsymbol\Omega\times\vect n'_k)\otimes\vect n'_k,
\]
两式逐项相同，所以 $\dd\vect n'_k/\dd t=\boldsymbol\Omega\times\vect n'_k$（第二个等号用了 $\sum_k\vect n'_k\otimes\vect n'_k=\unitd$）。

\textbf{无穷小情形对一下。}取 $S=\unitd+\dd{\boldsymbol\varphi}^{\times}$，则 $\dd S/\dd t=\dot{\boldsymbol\varphi}^{\times}$、$S^{\mathsf T}\approx\unitd$，所以
\[
\boldsymbol\Omega^{\times}=\dot{\boldsymbol\varphi}^{\times}
\quad\Longrightarrow\quad
\boldsymbol\Omega=\frac{\dd{\boldsymbol\varphi}}{\dd t}.
\]
角速度就是“转角矢量对时间的变化率”；方向是瞬时转轴，大小是转动的快慢。§5 的输运定理与欧拉方程都要用这个矢量。

\hdr{§4}{落点一：转动惯量张量}

到这里工具已经齐了：线性映射（§0）、二次型与主轴（§1）、并矢运算（§2）、旋转算子与角速度（§3）。从这一节起，每个结论都只用前面出现过的东西。

\subhdr{4.1\quad 从 $\vect L=\int\vect r\times\vect v\dd m$ 推出惯量张量（必读）}

取刚体绕定点 $O$ 转动（或取质心为基点），角速度为 $\boldsymbol\omega$。刚体上位置为 $\vect r$ 的点，其速度是
\[
\vect v=\boldsymbol\omega\times\vect r .
\]
为什么？绕某一瞬间的转轴转动的点，到轴的距离为 $\rho$ 时速度大小为 $\omega\rho$，方向垂直于轴与 $\vect r$ 张成的平面，正是 $\boldsymbol\omega\times\vect r$ 的方向；写成矢量就是上式。这是刚体运动学的基本结论（课件 p15 的总结也用了它），本稿从它出发。\aside{这里用到的两个词：\textbf{质心}是质量加权平均位置 $\vect R=\int\vect r\dd m/m$，§4.4 的平行轴定理会用到它的定义 $\int\vect r'\dd m=\vect0$；\textbf{瞬时轴}是某一瞬间速度为零的点的连线，角速度 $\boldsymbol\omega$ 就沿它。}

角动量（对 $O$ 点）定义为
\[
\vect L=\int\vect r\times\vect v\dd m
=\int\vect r\times(\boldsymbol\omega\times\vect r)\dd m .
\]
用三矢量叉乘公式 $\vect a\times(\vect b\times\vect c)=\vect b(\vect a\cdot\vect c)-\vect c(\vect a\cdot\vect b)$ 展开内层：
\[
\vect r\times(\boldsymbol\omega\times\vect r)
=\boldsymbol\omega(\vect r\cdot\vect r)-\vect r(\vect r\cdot\boldsymbol\omega)
=r^2\boldsymbol\omega-\vect r(\vect r\cdot\boldsymbol\omega).
\]
现在把两项都写成“某个张量作用在 $\boldsymbol\omega$ 上”：
\[
r^2\boldsymbol\omega=r^2\unitd\cdot\boldsymbol\omega,
\qquad
\vect r(\vect r\cdot\boldsymbol\omega)=(\vect r\otimes\vect r)\cdot\boldsymbol\omega .
\]
两个式子都作用在同一个 $\boldsymbol\omega$ 上，所以可以先把算子加起来再作用（这就是“$\ten J$ 是一台机器”的具体含义）。对 $\boldsymbol\omega$ 线性，所以可以提到积分外面：
\[
\boxed{\;\vect L=\ten J\cdot\boldsymbol\omega,\qquad
\ten J=\int\bigl(r^2\unitd-\vect r\otimes\vect r\bigr)\dd m\;}
\]
这就是\textbf{转动惯量张量}。它的结构值得读一遍：$r^2\unitd$ 让所有方向一视同仁，$\vect r\otimes\vect r$ 把沿 $\vect r$ 的那一份减掉。正因为有第二项，$\vect L$ 才一般\textbf{不平行}于 $\boldsymbol\omega$（§4.5 给具体数值）。

\textbf{离散核对。}对一个位于 $\vect r$、质量为 $m$ 的质点，
\[
\ten J=m(r^2\unitd-\vect r\otimes\vect r),
\qquad
\ten J\cdot\boldsymbol\omega=m\bigl[r^2\boldsymbol\omega-\vect r(\vect r\cdot\boldsymbol\omega)\bigr]=m\,\vect r\times(\boldsymbol\omega\times\vect r),
\]
与直接算 $m\vect r\times\vect v$ 一致。连续体就是把许多这样的质点加起来。

\subhdr{4.2\quad 分量、几何意义、为什么对称与正定（必读）}

把 $\ten J$ 逐分量写出来（用 $\delta_{ij}$ 表示单位并矢的分量）：
\[
J_{ij}=\int\bigl(r^2\delta_{ij}-x_ix_j\bigr)\dd m
\qquad\Longrightarrow\qquad
\ten J=
\begin{pmatrix}
\int(y^2+z^2)\dd m & -\int xy\dd m & -\int xz\dd m\\[3pt]
-\int xy\dd m & \int(x^2+z^2)\dd m & -\int yz\dd m\\[3pt]
-\int xz\dd m & -\int yz\dd m & \int(x^2+y^2)\dd m
\end{pmatrix}.
\]
读法：对角元是\textbf{绕对应坐标轴的转动惯量}；非对角元是\textbf{惯性积的负值}。为什么非对角元带负号？它就是 $-\vect r\otimes\vect r$ 那一项贡献的 $-x_ix_j$，没有别的理由。

\textbf{对称性是白送的}：$r^2\delta_{ij}-x_ix_j$ 对 $i,j$ 对称，所以 $J_{ij}=J_{ji}$。
\textbf{正定性也是白送的}（正定就是：对任意非零 $\boldsymbol\omega$，二次型 $\boldsymbol\omega\cdot\ten J\cdot\boldsymbol\omega>0$）：对任意 $\boldsymbol\omega$，
\[
\boldsymbol\omega\cdot\ten J\cdot\boldsymbol\omega
=\int\bigl[r^2\omega^2-(\vect r\cdot\boldsymbol\omega)^2\bigr]\dd m
=\int|\vect r\times\boldsymbol\omega|^2\dd m\ge0,
\]
因为 $|\vect r\times\boldsymbol\omega|^2=r^2\omega^2-(\vect r\cdot\boldsymbol\omega)^2$。等号只在“所有质量都在过原点且平行 $\boldsymbol\omega$ 的直线上”时成立；细杆沿自身方向时会出现这个退化（那根方向的转动惯量为零），一般三维刚体不会。所以惯量张量通常正定；但不管正定与否，它总是实对称，§1.5 的谱定理都可用：三个实特征值、三根正交主轴（正定只额外保证等值面是椭球）。

\textbf{动能}是同一个张量给出的二次型：
\[
K=\tfrac12\int v^2\dd m
=\tfrac12\int(\boldsymbol\omega\times\vect r)\cdot(\boldsymbol\omega\times\vect r)\dd m
=\tfrac12\boldsymbol\omega\cdot\ten J\cdot\boldsymbol\omega .
\]
这里出现了 §1.1 那句“课件 p9 的二次型”：$K$ 就是 $\boldsymbol\omega$ 的二次型，系数表就是 $\ten J$。

\subhdr{4.3\quad 主轴 = 惯量主轴，主值 = 主转动惯量（必读）}

在主轴系里 $\ten J=\operatorname{diag}(A,B,C)$（$A,B,C$ 就是三个\textbf{主转动惯量}），于是
\[
\vect L=(A\omega_1,\;B\omega_2,\;C\omega_3),
\qquad
K=\tfrac12\bigl(A\omega_1^2+B\omega_2^2+C\omega_3^2\bigr).
\]
三个主值满足三角形不等式
\[
A+B\ge C\quad\text{（以及另外两个循环）}.
\]
验证只需一行：
\[
A+B-C=\int\bigl[(y^2+z^2)+(x^2+z^2)-(x^2+y^2)\bigr]\dd m=2\int z^2\dd m\ge0 .
\]
这条不等式是很便宜的验算工具：算出来的三个主值若不满足，一定哪一步错了。

\textbf{对称轴的方向就是主轴。}如果刚体有一个几何对称轴，那么绕该轴转 $180^\circ$ 会把刚体映回自身，惯量张量在这个转动下不变。这个不变性迫使该轴与垂直平面之间的惯性积为零，于是该轴是一个特征方向；剩下两个主轴在垂直平面内由同样的谱定理给出（把 §1.5 的结论用在垂直平面内的 $2\times2$ 块上）。这条判断能省掉大量特征值计算：
\begin{itemize}
\item 均质长方体：三根棱方向都是主轴（每根轴转 $180^\circ$ 都不变）；
\item 正圆锥、圆柱：对称轴是主轴；另两根在垂直于它的平面内任意取（那两个主值相等）；
\item 均质正四面体：三根主轴过对棱中点连线，三个主值相等。
\end{itemize}

\begin{center}
\begin{tikzpicture}[scale=1.75,>={Stealth[length=1.7mm]}]
  \fill[bandviolet] (0,0)--(2.4,0)--(3.0,1.1)--(0.6,1.1)--cycle;
  \draw[violet,thick] (0,0)--(2.4,0)--(3.0,1.1)--(0.6,1.1)--cycle;
  \coordinate (C) at (1.5,0.55);
  \draw[vv] (0.45,0.55)--(2.85,0.55) node[right]{\scriptsize $J_1$};
  \draw[vv] (1.2845,0.1549)--(1.7155,0.9451) node[above right]{\scriptsize $J_2$};
  \draw[inkgray,thick] (C) circle (0.09);
  \fill (C) circle (0.02);
  \node[inkgray,above left] at (1.46,0.64) {\scriptsize $J_3$};
\end{tikzpicture}\\[2pt]
\figcap{\mbox{\textbf{图 5}（自绘）：}一块均质矩形\textbf{薄板}（厚度忽略，质量分布在板面内）的三根惯量主轴。$J_1$、$J_2$ 沿两条棱（在板面内），$J_3$ 垂直板面（画成 $\odot$）。斜投影下 $J_1$ 与 $J_2$ 看起来不垂直，实际在板面内是垂直的。绕 $J_3$、$J_1$ 或 $J_2$ 转 $180^\circ$，板都回到自己身上，所以三根都是主轴。垂直轴定理 $J_3=J_1+J_2$ \textbf{只对这样的平面薄板成立}（§4.4）。}
\end{center}

\subhdr{4.4\quad 常用结果与平行轴定理（必读）}

先算一个一维积分，后面两个公式都从它来。均质细杆长 $L$、质量 $m$，沿 $x$ 轴放，中点在原点。线密度是 $m/L$，所以
\[
\int_{-L/2}^{L/2}x^2\dd m
=\int_{-L/2}^{L/2}x^2\frac{m}{L}\dd x
=\frac{m}{L}\cdot\frac{x^3}{3}\Big|_{-L/2}^{L/2}
=\frac{mL^2}{12}.
\]
这就是细杆绕中点垂直轴的转动惯量；杆绕自身轴（$x$ 轴）的惯量为零，因为质量都在轴上。

\textbf{矩形薄板}（$a\times b$，板面为 $xy$ 面，厚度忽略）：
\[
J_x=\int y^2\dd m=\frac{mb^2}{12},\qquad
J_y=\int x^2\dd m=\frac{ma^2}{12},\qquad
J_z=\int(x^2+y^2)\dd m=\frac{m(a^2+b^2)}{12}.
\]
最后一条就是\textbf{垂直轴定理} $J_z=J_x+J_y$：$z$ 方向没有厚度，所以绕 $z$ 的惯量恰好是绕 $x$、$y$ 的惯量之和。它在并矢语言里不是新定理，只是一行读数。

\textbf{长方体}（边长 $a,b,c$，质量 $m$，质心为原点，棱沿坐标轴）：
\[
\ten J=\frac{m}{12}\operatorname{diag}\bigl(b^2+c^2,\;a^2+c^2,\;a^2+b^2\bigr).
\]
推导就是把 $\int(y^2+z^2)\dd m=\int y^2\dd m+\int z^2\dd m$ 分开，各用上面的杆公式；$y$ 方向长度 $b$ 给 $mb^2/12$，$z$ 方向长度 $c$ 给 $mc^2/12$。

\textbf{平行轴定理的并矢版。}把原点从质心 $C$ 挪到 $O$，设 $\vect R$ 是从 $O$ 指向 $C$ 的矢量，质元相对 $O$ 的位置是 $\vect r=\vect R+\vect r'$（$\vect r'$ 相对质心）。由质心定义 $\int\vect r'\dd m=\vect0$，展开
\[
\begin{aligned}
\ten J_O&=\int\bigl[(\vect R+\vect r')^2\unitd-(\vect R+\vect r')\otimes(\vect R+\vect r')\bigr]\dd m\\
&=\int\bigl(r'^2\unitd-\vect r'\otimes\vect r'\bigr)\dd m
+m\bigl(R^2\unitd-\vect R\otimes\vect R\bigr)\\
&\quad+\int\bigl[2(\vect R\cdot\vect r')\unitd-\vect R\otimes\vect r'-\vect r'\otimes\vect R\bigr]\dd m .
\end{aligned}
\]
最后一行三个积分的被积函数都含 $\vect r'$ 的一次因子，而 $\int\vect r'\dd m=\vect0$（质心定义），所以整行消失；剩下的第一项是 $\ten J_C$，第二项是 $m(R^2\unitd-\vect R\otimes\vect R)$。于是
\[
\ten J_O=\ten J_C+m\bigl(R^2\unitd-\vect R\otimes\vect R\bigr).
\]
\textbf{注意第二项是并矢，不能写成 $mR^2$ 一个标量}——绕不同方向的轴平移，增量不同。
标量版 $J=J_C+md^2$ 只是它在“轴平行且过 $O$”这一个特例上的读数：取单位轴 $\vect n$，则
\[
\vect n\cdot\bigl[m(R^2\unitd-\vect R\otimes\vect R)\bigr]\cdot\vect n
=m\bigl(R^2-(\vect n\cdot\vect R)^2\bigr)=m d^2,
\]
其中 $d$ 是质心到轴的垂直距离。

\warnbox{\textbf{选读（可跳过）：椭球的惯量张量。}
这一段给一个更复杂的例子，回答“非长方体怎么算”，不影响后面；跳过后从 §4.5 继续。
\par\smallskip
均质椭球 $x^2/a^2+y^2/b^2+z^2/c^2\le1$，质量 $M$，密度 $\rho$，体积 $V=4\pi abc/3$。用伸缩坐标 $x=a\xi,\;y=b\eta,\;z=c\zeta$，椭球变成单位球，$\dd V=abc\,\dd V_{\text{球}}$。于是
\[
\int_E y^2\dd V=abc\,b^2\int_{\text{球}}\eta^2\dd V_{\text{球}}
=abc\,b^2\cdot\frac{4\pi}{15}
=\frac{4\pi abc}{3}\cdot\frac{b^2}{5}=V\frac{b^2}{5},
\]
其中单位球的 $\int\eta^2\dd V_{\text{球}}=4\pi/15$ 可由球坐标算得（$\int_0^1r^4\dd r=\tfrac15$、$\int_0^\pi\sin^3\theta\dd\theta=\tfrac43$、$\int_0^{2\pi}\cos^2\varphi\dd\varphi=\pi$，相乘得 $4\pi/15$）。同理 $\int x^2\dd V=Va^2/5$、$\int z^2\dd V=Vc^2/5$。所以
\[
\ten J=\frac{M}{5}\operatorname{diag}\bigl(b^2+c^2,\;a^2+c^2,\;a^2+b^2\bigr).
\]
球 $a=b=c=R$ 时 $\ten J=\tfrac25MR^2\unitd$，各向同性，任何方向都是主轴。这个结果与 RIT Whelan notes 10 的 pp.11--12 一致（那里用同一组伸缩坐标算的）。}

\subhdr{4.5\quad 教学示范：$\vect L$ 到底偏多少}

取均质长方体 $a=1$、$b=2$、$c=3$ m，质量 $m=12$ kg，质心在原点。由 §4.4，
\[
\ten J=\frac{12}{12}\operatorname{diag}(2^2+3^2,\;1^2+3^2,\;1^2+2^2)
=\operatorname{diag}(13,\;10,\;5)\ \text{kg·m}^2 .
\]
让它绕一个与三条棱都成等角的方向转，$\boldsymbol\omega=\omega_0(1,1,1)/\sqrt3$，取 $\omega_0=1$ rad/s。（注意：这\textbf{不是}这个长方体的空间对角线；从 $(0,0,0)$ 到 $(1,2,3)$ 的空间对角线方向是 $(1,2,3)/\sqrt{14}$。这里选 $(1,1,1)/\sqrt3$，是因为它与三条棱的夹角相等、计算干净，而下面的结论只用到“$\boldsymbol\omega$ 与三根主轴都不平行”。）则
\[
\vect L=\ten J\cdot\boldsymbol\omega=\frac{1}{\sqrt3}(13,\;10,\;5)\approx(7.506,\;5.774,\;2.887)\ \text{kg·m}^2/\text{s}.
\]
$\vect L$ 与 $\boldsymbol\omega\propto(1,1,1)$ 明显不平行。夹角由
\[
\cos\theta=\frac{\vect L\cdot\boldsymbol\omega}{|\vect L|\,|\boldsymbol\omega|}
=\frac{28/3}{\sqrt{294/3}}
=\frac{28}{3\sqrt{98}}=\frac{4}{3\sqrt2}=\frac{2\sqrt2}{3}\approx0.9428
\quad\Longrightarrow\quad
\theta\approx19.47^\circ .
\]
核对：$|\vect L|=\sqrt{13^2+10^2+5^2}/\sqrt3=\sqrt{294/3}=\sqrt{98}=7\sqrt2$，$\vect L\cdot\boldsymbol\omega=(13+10+5)/3=28/3$，两者相除时 $\sqrt3$ 与公因子都约掉；$\sin\theta=\sqrt{1-8/9}=1/3$。

\textbf{若强行让它不晃，需要多大力矩？}若 $\boldsymbol\omega$ 在体坐标系里保持不动，则 $\vect L$ 在体坐标系里也不动，$\dd\vect L/\dd t|_{\text{体}}=\vect0$，固定系里的变化只剩“整体转动”那一项（§5.1）：
\[
\vect M=\boldsymbol\omega\times\vect L
=\frac{1}{3}\begin{pmatrix}1\\1\\1\end{pmatrix}\times\begin{pmatrix}13\\10\\5\end{pmatrix}
=\frac{1}{3}\begin{pmatrix}5-10\\13-5\\10-13\end{pmatrix}
=\frac{1}{3}\begin{pmatrix}-5\\8\\-3\end{pmatrix},
\]
\[
|\vect M|=\frac{1}{3}\sqrt{25+64+9}=\frac{\sqrt{98}}{3}=\frac{7\sqrt2}{3}\approx3.30\ \text{N·m}.
\]
方向垂直于 $\boldsymbol\omega$ 与 $\vect L$ 张成的平面（图上就是“把 $\vect L$ 往 $\boldsymbol\omega$ 那边扳”的方向）。\textbf{自由转动时没有这个力矩，所以 $\boldsymbol\omega$ 会自己动起来}——这就是“刚体转起来会晃”的数学原因，§5.4 会把这个力矩与欧拉方程接上。

\hdr{§5}{落点二：欧拉方程与对称陀螺}

\subhdr{5.1\quad 输运定理：为什么不能直接对 $\vect L$ 求导（必读）}

初学者最常写的一步错是
\[
\frac{\dd\vect L}{\dd t}=\ten J\cdot\frac{\dd\boldsymbol\omega}{\dd t}
\qquad\text{（错，除非 }\ten J\text{ 在该系里是常量）}.
\]
$\ten J$ 是“跟着刚体转”的张量。在固定系里看，它的九个分量随时间变；在体坐标系里看，它才是不变的。两个系里的时间导数由一个公式联系，叫\textbf{输运定理}。

\begin{center}
\begin{tikzpicture}[scale=1.10,>={Stealth[length=1.7mm]}]
  % 左：固定系
  \begin{scope}
    \draw[ax] (-1.5,0)--(1.65,0) node[right]{\scriptsize $\vect e_1$};
    \draw[ax] (0,-1.15)--(0,1.6) node[above]{\scriptsize $\vect e_2$};
    \begin{scope}[rotate=32]
      \draw[violet,thick,fill=bandviolet] (0,0) ellipse (1.05 and 0.42);
    \end{scope}
    \draw[violet,densely dashed,->] (-1.06,-0.66)--(1.06,0.66);
    \draw[violet,densely dashed,->] (0.53,-0.85)--(-0.53,0.85);
    \node[violet] at (1.20,0.80) {\scriptsize $\vect n'_1$};
    \node[violet] at (-0.72,0.92) {\scriptsize $\vect n'_2$};
    \draw[inkgray,->] (0.34,1.16) arc (72:38:1.18);
    \node[inkgray] at (0.92,1.36) {\scriptsize $\boldsymbol\Omega$};
    \node[align=center] at (0,-1.70) {\scriptsize 固定系：$\vect n'_k(t)$ 在转，\\ $\ten J$ 的九个数随时间变};
  \end{scope}
  % 右：体系
  \begin{scope}[xshift=6.4cm]
    \draw[ax] (-1.36,0.848)--(1.62,-1.012) node[above right]{\scriptsize $\vect e_1$};
    \draw[ax] (-0.795,-1.272)--(0.901,1.442) node[above right]{\scriptsize $\vect e_2$};
    \draw[violet,thick,fill=bandviolet] (0,0) ellipse (1.05 and 0.42);
    \draw[violet,thick,->] (0,0)--(1.45,0) node[above]{\scriptsize $\vect n'_1$};
    \draw[violet,thick,->] (0,0)--(0,1.25) node[above right]{\scriptsize $\vect n'_2$};
    \node[align=center] at (0,-1.70) {\scriptsize 体系：$\vect n'_k$ 定死（画面上是正的），\\ $\ten J=\operatorname{diag}(A,B,C)$ 不变};
  \end{scope}
\end{tikzpicture}\\[2pt]
\figcap{\mbox{\textbf{图 6}（自绘）：}同一块刚体，两种看法。左边在固定系里看：主轴 $\vect n'_k$ 在转，所以 $\ten J$ 的账本随时间变，“$\dd\vect L/\dd t=\ten J\cdot\dd\boldsymbol\omega/\dd t$”是错的。右边改用体系：把 $\vect n'_k$ 当成不动的坐标轴（灰虚线的 $\vect e_k$ 反过来在转），$\ten J$ 的账本就是常量 $\operatorname{diag}(A,B,C)$，代价是多出与转动有关的项（陀螺项）。}
\end{center}

\textbf{矢量版。}设随体基 $\vect n'_k$ 以角速度 $\boldsymbol\Omega$ 转动（§3.5 的结论：$\dd\vect n'_k/\dd t=\boldsymbol\Omega\times\vect n'_k$）。把矢量写成 $\vect A=A_k\vect n'_k$，其中 $A_k$ 是体坐标分量。求导：
\[
\Bigl(\frac{\dd\vect A}{\dd t}\Bigr)_{\text{固定}}
=\dot A_k\vect n'_k+A_k\frac{\dd\vect n'_k}{\dd t}
=\Bigl(\frac{\dd\vect A}{\dd t}\Bigr)_{\text{体}}+\boldsymbol\Omega\times\vect A .
\]
第一项是“体坐标分量变化”的贡献，第二项是“基矢量自己在转”的贡献。这就是相对运动里熟悉的输运公式。

\textbf{并矢版。}设 $\ten D=D_{ij}\vect n'_i\otimes\vect n'_j$（$D_{ij}$ 是体坐标分量）。逐项求导时三项都不能漏：
\[
\begin{aligned}
\Bigl(\frac{\dd\ten D}{\dd t}\Bigr)_{\text{固定}}
&=\dot D_{ij}\,\vect n'_i\otimes\vect n'_j
+D_{ij}\Bigl(\frac{\dd\vect n'_i}{\dd t}\Bigr)\otimes\vect n'_j
+D_{ij}\,\vect n'_i\otimes\Bigl(\frac{\dd\vect n'_j}{\dd t}\Bigr)\\
&=\Bigl(\frac{\dd\ten D}{\dd t}\Bigr)_{\text{体}}
+D_{ij}(\boldsymbol\Omega\times\vect n'_i)\otimes\vect n'_j
+D_{ij}\,\vect n'_i\otimes(\boldsymbol\Omega\times\vect n'_j).
\end{aligned}
\]
第一项是体分量自己的变化率，就是 $(\dd\ten D/\dd t)_{\text{体}}$。后两项各把一个基矢量换掉；为把“换基”写成一台机器，用 §2.6 的对偶张量 $\boldsymbol\Omega^{\times}$（它就是矩阵乘法意义下的反对称矩阵），并规定它对并矢的作用为
\[
\boldsymbol\Omega^{\times}\cdot(\vect a\otimes\vect b)\stackrel{\text{def}}{=}(\boldsymbol\Omega\times\vect a)\otimes\vect b,
\qquad
(\vect a\otimes\vect b)\cdot\boldsymbol\Omega^{\times}\stackrel{\text{def}}{=}\vect a\otimes(\vect b\cdot\boldsymbol\Omega^{\times})
=-\vect a\otimes(\boldsymbol\Omega\times\vect b).
\]
第二个定义用了 $(\boldsymbol\Omega^{\times})^{\mathsf T}=-\boldsymbol\Omega^{\times}$，所以右边的负号不是另加的。按线性，后两项正好是 $\boldsymbol\Omega^{\times}\cdot\ten D$ 与 $-\ten D\cdot\boldsymbol\Omega^{\times}$，于是
\[
\boxed{\;\Bigl(\frac{\dd\ten D}{\dd t}\Bigr)_{\text{固定}}
=\Bigl(\frac{\dd\ten D}{\dd t}\Bigr)_{\text{体}}
+\boldsymbol\Omega^{\times}\cdot\ten D-\ten D\cdot\boldsymbol\Omega^{\times}\;}
\]
这里两个乘积都是\textbf{普通矩阵乘法}：$\boldsymbol\Omega^{\times}\cdot\ten D$ 和 $\ten D\cdot\boldsymbol\Omega^{\times}$ 各把张量的一个指标“转一次”，两个指标地位相反，所以一正一负。

\textbf{分量写法。}按 §2.6 的矩阵定义，两个乘积的分量就是
\[
(\boldsymbol\Omega^{\times}\cdot\ten D)_{ij}
=\sum_k(\boldsymbol\Omega^{\times})_{ik}D_{kj},
\qquad
(\ten D\cdot\boldsymbol\Omega^{\times})_{ij}
=\sum_kD_{ik}(\boldsymbol\Omega^{\times})_{kj}.
\]
若再用 §2.6 的 $\varepsilon$ 写法 $(\boldsymbol\Omega^{\times})_{ik}=\varepsilon_{imk}\Omega_m$，就是
\[
(\boldsymbol\Omega^{\times}\cdot\ten D)_{ij}
=\sum_{k,m}\varepsilon_{imk}\Omega_mD_{kj},
\qquad
(\ten D\cdot\boldsymbol\Omega^{\times})_{ij}
=\sum_{k,m}D_{ik}\varepsilon_{kmj}\Omega_m .
\]
这两个式子与上面的并矢定义逐项一致，没有额外负号。

\textbf{单并矢核对。}取 $\ten D=\vect a\otimes\vect b$，按定义
\[
\boldsymbol\Omega^{\times}\cdot\ten D=(\boldsymbol\Omega\times\vect a)\otimes\vect b,
\qquad
\ten D\cdot\boldsymbol\Omega^{\times}=-\vect a\otimes(\boldsymbol\Omega\times\vect b).
\]
所以右边两项合起来是
\[
(\boldsymbol\Omega\times\vect a)\otimes\vect b+\vect a\otimes(\boldsymbol\Omega\times\vect b),
\]
正好是乘积求导的两项。用 $\vect a=\vect e_1$、$\vect b=\vect e_2$、$\boldsymbol\Omega=\vect e_3$ 核验：
\[
(\boldsymbol\Omega\times\vect e_1)\otimes\vect e_2+\vect e_1\otimes(\boldsymbol\Omega\times\vect e_2)
=\vect e_2\otimes\vect e_2-\vect e_1\otimes\vect e_1
=\operatorname{diag}(-1,1,0),
\]
而
\[
\boldsymbol\Omega^{\times}\cdot(\vect e_1\otimes\vect e_2)-(\vect e_1\otimes\vect e_2)\cdot\boldsymbol\Omega^{\times}
=\vect e_2\otimes\vect e_2-\vect e_1\otimes\vect e_1,
\]
两边相同。这个例子既非零又非对称，能暴露转置或符号错误。

\textbf{单位张量的检查。}取 $\ten D=\unitd$，左边为 $\vect0$（单位张量永远是单位张量），右边
\[
\boldsymbol\Omega^{\times}\cdot\unitd-\unitd\cdot\boldsymbol\Omega^{\times}
=\boldsymbol\Omega^{\times}-\boldsymbol\Omega^{\times}=\vect0 .
\]

\subhdr{5.2\quad 体坐标系里的角动量定理与欧拉方程（必读）}

固定系里角动量定理是
\[
\Bigl(\frac{\dd\vect L}{\dd t}\Bigr)_{\text{固定}}=\vect M,
\]
其中 $\vect M$ 是外力对定点（或质心）的力矩。力矩就是“力对点的矩” $\vect M=\int\vect r\times\dd\vect F$；这条角动量定理是牛顿第二定律对质点系求和的结果，条件是惯性系、对定点或质心取矩。
换成体坐标系（体坐标系以 $\boldsymbol\omega$ 转动），用矢量版输运公式：
\[
\Bigl(\frac{\dd\vect L}{\dd t}\Bigr)_{\text{体}}+\boldsymbol\omega\times\vect L=\vect M .
\]
体坐标系里 $\ten J=\operatorname{diag}(A,B,C)$ 是常量，$\vect L=(A\omega_1,B\omega_2,C\omega_3)$，于是 $(\dd\vect L/\dd t)_{\text{体}}=(A\dot\omega_1,B\dot\omega_2,C\dot\omega_3)$。再算 $\boldsymbol\omega\times\vect L$ 的三个分量：
\[
(\boldsymbol\omega\times\vect L)_1=\omega_2(C\omega_3)-\omega_3(B\omega_2)=(C-B)\omega_2\omega_3,
\]
另两个循环。合起来就是\textbf{欧拉方程}：
\[
\boxed{
\begin{aligned}
A\dot\omega_1+(C-B)\omega_2\omega_3&=M_1,\\
B\dot\omega_2+(A-C)\omega_3\omega_1&=M_2,\\
C\dot\omega_3+(B-A)\omega_1\omega_2&=M_3 .
\end{aligned}}
\]
无外力矩（$M_k=0$）时，这三条就是刚体自由转动的全部动力学。$A,B,C$ 是体坐标系里的三个主转动惯量，$\omega_k$ 是角速度在体坐标系的三个分量，$M_k$ 是外力矩的体坐标分量。

\subhdr{5.3\quad 对称陀螺：$A=B\ne C$（必读）}

取 $A=B$、$\vect M=\vect0$。第三条方程立刻给 $C\dot\omega_3=0$，即
\[
\omega_3=\text{常量}.
\]
前两条变成
\[
\dot\omega_1=-\Omega\,\omega_2,\qquad \dot\omega_2=\Omega\,\omega_1,
\qquad
\Omega\equiv\frac{C-A}{A}\,\omega_3 .
\]
这两条不借助微分方程理论也能读懂。把横向部分看成一个二维矢量 $\vect w_\perp=(\omega_1,\omega_2)$，上式说
\[
\frac{\dd}{\dd t}\begin{pmatrix}\omega_1\\\omega_2\end{pmatrix}
=\Omega\begin{pmatrix}-\omega_2\\\omega_1\end{pmatrix},
\]
即横向“速度”与 $\vect w_\perp$ 垂直、大小是 $|\Omega|\,|\vect w_\perp|$。先看大小：把上式与 $(\omega_1,\omega_2)$ 点乘，
\[
\omega_1\dot\omega_1+\omega_2\dot\omega_2
=\Omega(-\omega_1\omega_2+\omega_2\omega_1)=0,
\]
所以 $\omega_1^2+\omega_2^2=\omega_\perp^2$ 是常数。再看方向：若 $\Omega>0$，速度总是把 $(\omega_1,\omega_2)$ 朝逆时针方向转 $90^\circ$；所以它在以原点为圆心、半径 $\omega_\perp$ 的圆上匀速转，角速度就是 $\Omega$（$\Omega<0$ 时反向）。

把这句话写成公式：
\[
\omega_1=\omega_\perp\cos(\Omega t+\varphi_0),\qquad
\omega_2=\omega_\perp\sin(\Omega t+\varphi_0),
\]
其中 $\omega_\perp\ge0$ 是圆半径，$\varphi_0$ 是 $t=0$ 时横向矢量与 $\vect e_1$ 的夹角。直接代回可以当场核对：
\[
\dot\omega_1=-\Omega\omega_\perp\sin(\Omega t+\varphi_0)=-\Omega\omega_2,
\qquad
\dot\omega_2=\Omega\omega_\perp\cos(\Omega t+\varphi_0)=\Omega\omega_1 .
\]
给定初始横向矢量 $\vect w_\perp(0)=(\omega_1(0),\omega_2(0))$，两个常数就定下来了（$\omega_\perp>0$ 时；若 $\omega_\perp=0$，横向分量恒为零，$\varphi_0$ 可任取）：
\[
\omega_\perp=\sqrt{\omega_1(0)^2+\omega_2(0)^2},
\qquad
\cos\varphi_0=\frac{\omega_1(0)}{\omega_\perp},\quad
\sin\varphi_0=\frac{\omega_2(0)}{\omega_\perp}.
\]
例如 $t=0$ 时 $\vect w_\perp$ 沿 $\vect e_1$，则 $\varphi_0=0$，解是 $\omega_1=\omega_\perp\cos\Omega t$、$\omega_2=\omega_\perp\sin\Omega t$；若沿 $\vect e_2$，则 $\varphi_0=\pi/2$。这样从任意给定的初始横向矢量都能写出解、也能代回核对，不需要预先学过“简谐振子”。（同一个数学结构在其它问题里常叫“简谐振子”；本稿不需要那个名称，也不依赖它的结论。）
\textbf{结论：}角速度矢量的横向分量在体坐标系里绕体轴 $3$ 匀速转动，角速度 $\Omega$（可正可负，负号表示反向）、大小 $\omega_\perp$ 不变，$\omega_3$ 也不变，所以 $|\boldsymbol\omega|$ 不变。
（严格说，这里是 $\boldsymbol\omega$ 在体坐标系里绕 $\vect e_3$ 转，不是陀螺本身在空间中的进动；后者还要再换回固定系。所谓\textbf{进动}，就是转轴或角速度方向绕着另一条轴慢慢转。）

\textbf{几何图像。}$\vect L=(A\omega_1,A\omega_2,C\omega_3)$。因为 $\vect L$ 的前两个分量只是 $\boldsymbol\omega$ 的前两个分量乘同一个常数 $A$，所以 $\boldsymbol\omega$、$\vect L$、$\vect e_3$ \textbf{三个矢量始终共面}；这个平面绕 $\vect e_3$ 以 $\Omega$ 转。这是对称陀螺最有用的图像。

\begin{center}
\begin{tikzpicture}[scale=1.6,>={Stealth[length=1.7mm]}]
  \coordinate (O) at (0,0);
  \coordinate (E3) at (0,1.6);
  \coordinate (W) at (0.937,1.040);
  \coordinate (L) at (0.876,1.217);
  \draw[ax] (O)--(E3) node[above]{\scriptsize $\vect e_3$};
  \draw[inkgray,thin,densely dashed] (0.937,1.040)--(0.937,0);
  \draw[inkgray,thin,densely dashed] (0.876,1.217)--(0,1.217);
  \draw[vv] (O)--(W) node[right]{\scriptsize $\boldsymbol\omega$};
  \draw[violet,very thick,->,>={Stealth[length=1.7mm]}] (O)--(L) node[above left]{\scriptsize $\vect L$};
  \draw[inkgray,->] (0.30,1.42) arc (20:160:0.30);
  \node[inkgray] at (0.44,1.52) {\scriptsize $\Omega$};
  \fill (O) circle (0.7pt);
\end{tikzpicture}\\[2pt]
\figcap{\mbox{\textbf{图 7}（自绘）：}对称陀螺的 $\boldsymbol\omega$、$\vect L$、$\vect e_3$ 共面。图中 $A<C$（扁盘状）的取法使 $\vect L$ 比 $\boldsymbol\omega$ 更靠近 $\vect e_3$；这个平面绕 $\vect e_3$ 以 $\Omega$ 转，$\boldsymbol\omega$ 的横向分量在体坐标系里画出一个圆。}
\end{center}

\begin{center}
\fcolorbox{violet}{bandviolet}{\parbox{0.94\linewidth}{\centering\vspace{3pt}
\small $C>A$（扁盘状，如地球）时 $\Omega>0$，$\boldsymbol\omega$ 的横向分量与 $\vect L$ 绕 $\vect e_3$ 同向转；$C<A$（细长杆状）时反向。\vspace{3pt}}}
\end{center}

\textbf{教学示范。}取对称陀螺 $A=B=4$ kg·m$^2$、$C=5$ kg·m$^2$，自转 $\omega_3=6$ rad/s，无外力矩。则
\[
\Omega=\frac{C-A}{A}\omega_3=\frac{5-4}{4}\times6=1.5\ \text{rad/s},
\qquad
T=\frac{2\pi}{\Omega}=\frac{2\pi}{1.5}=\frac{4\pi}{3}\approx4.19\ \text{s}.
\]
即 $\boldsymbol\omega$ 的横向分量每 $4.19$ 秒绕体轴转一圈，$\omega_3$ 与 $|\boldsymbol\omega|$ 都不变。$C>A$，属于“扁盘状”，$\Omega>0$，横向分量沿 $+\vect e_3$ 的正向转。

\subhdr{5.4\quad 回到 §4.5：维持对角转动需要多大力矩（必读）}

把“$\vect L$ 与 $\boldsymbol\omega$ 不平行”代进 $\vect M=\boldsymbol\omega\times\vect L$，就得到“维持这种转动需要多大扭矩”。§4.5 已经对 $m=12$、$a=1,b=2,c=3$ 的长方体算过：$|\vect M|=\sqrt{98}/3\approx3.30$ N·m，方向沿 $(1/3)(-5,8,-3)$。

一般地，若 $\boldsymbol\omega$ 在体坐标系里保持不动，则
\[
\vect M=\boldsymbol\omega\times\vect L,\qquad
|\vect M|=|\boldsymbol\omega|\,|\vect L|\sin\theta,
\]
其中 $\theta$ 是 $\boldsymbol\omega$ 与 $\vect L$ 的夹角；方向垂直于 $\boldsymbol\omega$ 与 $\vect L$ 张成的平面。
\textbf{自由转动时没有这个力矩，所以 $\boldsymbol\omega$ 必须自己动起来}——这正是 §5.3 那条进动的来源。反过来说：\textbf{若 $\vect M=\vect0$ 且转动是稳定的，则 $\boldsymbol\omega\times\vect L=\vect0$，即 $\vect L\parallel\boldsymbol\omega$，$\boldsymbol\omega$ 必须沿某根主轴。}这就是“只有绕主轴的匀速转动可以无扭矩维持”的并矢解释。

\hdr{§6}{落点三：应力张量}

\subhdr{6.1\quad 面力是法向的函数：柯西应力与商法则（必读）}

在连续体里（把物体看成连续分布的物质，而不是一个个离散质点）过一点取一个截面，法向为单位矢量 $\vect n$。截面上单位面积受到的力记作 $\vect t(\vect n)$，叫\textbf{面力矢量}（traction）。它描述的是“法向为 $\vect n$ 的那一侧材料，对另一侧单位面积的作用力”。

\textbf{柯西应力定理说：$\vect t$ 是 $\vect n$ 的线性函数。}为什么？取一个微小的四面体：三个面分别垂直于 $\vect e_1,\vect e_2,\vect e_3$，第四个面的法向是 $\vect n$。设第四个面的面积为 $\dd S$，三个坐标面的面积分别是 $\dd S\,n_1,\dd S\,n_2,\dd S\,n_3$（面积投影：一个面的面积乘上它与坐标面的夹角的余弦，正好是法向分量）。四面体静止时合力为零；体积力正比于体积（三阶小量），当四面体缩小时可以忽略，只剩下四个面上的面力：
\[
\vect t(\vect n)\dd S+\vect t(-\vect e_1)\dd S\,n_1+\vect t(-\vect e_2)\dd S\,n_2+\vect t(-\vect e_3)\dd S\,n_3=\vect0 .
\]
用 $\vect t(-\vect e_j)=-\vect t(\vect e_j)$（作用与反作用），除 $\dd S$ 得
\[
\vect t(\vect n)=n_1\vect t(\vect e_1)+n_2\vect t(\vect e_2)+n_3\vect t(\vect e_3).
\]
\textbf{这就是线性性，而且是“对任意 $\vect n$ 都成立”的线性性。}把 $\vect t(\vect e_j)$ 的三个分量记作 $\sigma_{ij}$：
\[
\vect t(\vect e_j)=\sigma_{ij}\vect e_i,
\qquad
\boxed{\;\vect t(\vect n)=\ten\sigma\cdot\vect n,\qquad t_i=\sigma_{ij}n_j\;}
\]
其中 $\sigma_{ij}$ 读作“\textbf{法向为 $\vect e_j$ 的面上、沿 $\vect e_i$ 方向的应力分量}”。第一个指标是力的方向，第二个指标是面的法向——这个顺序在 §6.4 的散度里会再出现一次。

\textbf{为什么它一定是张量？}因为 $\vect t$ 与 $\vect n$ 都是矢量，而“对任意 $\vect n$ 成立 $\vect t=\ten\sigma\vect n$”正是 §2.7 商法则的前提，所以 $\ten\sigma$ 自动是二阶张量。\textbf{这一步不是假设，是被逼出来的}——这也是 §2.7 存在的理由。

\subhdr{6.2\quad 对称性来自角动量守恒（必读）}

对微元取力矩平衡（不设体力矩、准静态；准静态就是每个时刻都可当作平衡处理，加速度的影响可忽略）。考虑一个小立方体：作用在法向为 $\vect e_1$ 的面上的切应力 $\sigma_{21}$，与作用在法向为 $\vect e_2$ 的面上的切应力 $\sigma_{12}$，分别产生绕第三个轴的力矩。两个力臂相同、面积相同；它们产生的力矩方向相反，平衡时必须等大，于是两个切应力分量相等：
\[
\sigma_{ij}=\sigma_{ji}.
\]
三个独立等式给出 $\sigma_{12}=\sigma_{21}$、$\sigma_{13}=\sigma_{31}$、$\sigma_{23}=\sigma_{32}$。所以应力张量也\textbf{对称}，§1.5 的谱定理同样适用：三个\textbf{主应力} $\sigma_1,\sigma_2,\sigma_3$、三个互相垂直的\textbf{主平面}（这些截面上剪应力为零）。主应力就是 $\ten\sigma$ 的特征值，主平面的法向就是特征向量。

\textbf{一个与坐标无关的数。}三个主应力之和
\[
\sigma_1+\sigma_2+\sigma_3=\tr\ten\sigma
\]
是与坐标无关的（§0.5）；它对应“静水压强”那种不随截面变化的成分，$p=-\tfrac13\tr\ten\sigma$（拉伸为正时，压力取负号）。把各向同性的压力部分减掉，剩下的部分叫偏斜部分，它才决定材料怎么变形（这一段不需要计算，知道有这层区分即可）。

\subhdr{6.3\quad 教学示范：斜截面上的面力分解（必读）}

取单轴拉伸
\[
\ten\sigma=\sigma_0\,\vect e_1\otimes\vect e_1
\qquad(\sigma_0>0),
\]
即只有 $\sigma_{11}=\sigma_0$，其余分量为零。主轴就是 $\vect e_1,\vect e_2,\vect e_3$，主应力是 $\sigma_0,0,0$。

在 $xy$ 面内取与 $\vect e_1$ 成 $\theta=30^\circ$ 的截面，法向
\[
\vect n=\cos\theta\,\vect e_1+\sin\theta\,\vect e_2=\frac{\sqrt3}{2}\vect e_1+\frac12\vect e_2 .
\]
面力：
\[
\vect t=\ten\sigma\cdot\vect n
=\sigma_0\,\vect e_1(\vect e_1\cdot\vect n)
=\sigma_0\cos\theta\,\vect e_1
=\frac{\sqrt3}{2}\sigma_0\,\vect e_1 .
\]
它仍然沿 $\vect e_1$，但斜截面把它切成两个分量。法向分量：
\[
t_n=\vect t\cdot\vect n=\sigma_0\cos^2\theta=\frac34\sigma_0,
\]
法向矢量是 $t_n\vect n$。切向矢量是 $\vect t-t_n\vect n$，其大小
\[
\tau=\sigma_0\sin\theta\cos\theta=\frac{\sqrt3}{4}\sigma_0,
\]
方向是 $(\sin\theta,-\cos\theta)=(1/2,-\sqrt3/2)$（沿截面方向）。核对：$|\vect t|^2=\tfrac34\sigma_0^2$，$t_n^2+\tau^2=\tfrac{9}{16}\sigma_0^2+\tfrac{3}{16}\sigma_0^2=\tfrac34\sigma_0^2$，一致。

\begin{center}
\begin{tikzpicture}[scale=1.65,>={Stealth[length=1.7mm]}]
  % 左：试件与斜截面
  \begin{scope}
    \fill[bandcream] (0,-0.6) rectangle (2.4,0.6);
    \draw[violet,thick] (0,-0.6) rectangle (2.4,0.6);
    \draw[very thick,violet!60!black] (0.9,0.5196)--(1.5,-0.5196);
    \draw[vv] (1.2,0)--(1.7369,0.31) node[above right]{\scriptsize $\vect n$};
    \draw[vv] (1.2,0)--(2.25,0) node[below right]{\scriptsize $\vect t=\ten\sigma\cdot\vect n$};
    \draw[vv] (2.4,0)--(3.0,0) node[right]{\scriptsize $+\sigma_0\vect e_1$};
    \draw[vv] (0,0)--(-0.6,0) node[left]{\scriptsize $-\sigma_0\vect e_1$};
    \node at (1.2,-0.95) {\scriptsize 单轴拉伸：两端向外拉};
  \end{scope}
  % 右：t 沿 n（30°）与截面方向分解
  \begin{scope}[xshift=6.1cm]
    \draw[ax] (-0.26,-0.15)--(1.55,0.895) node[above right]{\scriptsize $\vect n$};
    \draw[vv] (0,0)--(1.55,0) node[below right]{\scriptsize $\vect t$};
    \draw[vviolet] (0,0)--(1.1626,0.6712);
    \draw[inkgray,very thick,->,>={Stealth[length=1.7mm]}] (0,0)--(0.3875,-0.6712);
    \draw[inkgray,densely dotted] (1.1626,0.6712)--(1.55,0);
    \draw[inkgray,densely dotted] (0.3875,-0.6712)--(1.55,0);
    \draw[inkgray,thin] (0.1039,0.0600)--(0.1639,-0.0439)--(0.0600,-0.1039);
    \node[violet,above left] at (1.08,0.70) {\scriptsize $\tfrac34\sigma_0$};
    \node[inkgray,below right] at (0.36,-0.72) {\scriptsize $\tfrac{\sqrt3}{4}\sigma_0$};
    \draw[->] (0.45,0) arc (0:30:0.45);
    \node at (0.60,0.20) {\scriptsize $30^\circ$};
  \end{scope}
\end{tikzpicture}\\[2pt]
\figcap{\mbox{\textbf{图 8}（自绘，示意；各向量长度不按同一比例）：}单轴拉伸 $\ten\sigma=\sigma_0\vect e_1\otimes\vect e_1$。左：右端面的外法向是 $+\vect e_1$，面力 $\vect t(+\vect e_1)=\ten\sigma\cdot\vect e_1=+\sigma_0\vect e_1$；左端面的外法向是 $-\vect e_1$，面力 $\vect t(-\vect e_1)=\ten\sigma\cdot(-\vect e_1)=-\sigma_0\vect e_1$。两个向量都指向试件外侧，所以是“拉”而不是“压”。$30^\circ$ 截面（法向 $\vect n$）上的面力 $\vect t=\ten\sigma\cdot\vect n$ 仍沿 $+\vect e_1$。右：把这个 $\vect t$ 分成法向 $\tfrac34\sigma_0$ 与切向 $\tfrac{\sqrt3}{4}\sigma_0$ 两个分量；灰色虚线是 $\vect n$ 方向的参考轴（$30^\circ$），不是原来的 $x$ 轴。\textbf{斜面上原来看不出有剪应力，换个截面就出来了}——这就是 $\ten\sigma\cdot\vect n$ 的全部内容。最大剪应力出现在哪个角度、等于多少，是练习 P10 的内容。}
\end{center}

\subhdr{6.4\quad 平衡方程：应力变化怎样变成净力（必读）}

\textbf{偏导在做什么。}应力张量在物体里逐点不同，是一个\textbf{张量场} $\ten\sigma(\vect x)$：每个点给一个 $3\times3$ 表。$\partial_j=\partial/\partial x_j$ 的意思是：\textbf{只让第 $j$ 个坐标变，另外两个坐标固定}，看分量 $\sigma_{ij}$ 随它怎么变。例如 $\partial_1\sigma_{11}$ 是“沿 $x_1$ 方向走一小步，$\sigma_{11}$ 增加多少除以这一步”，$x_2,x_3$ 不动。

\textbf{为什么它的和是净面力。}取一个边长为 $h$ 的小立方体，考虑法向为 $\pm\vect e_j$ 的一对面上、沿 $\vect e_i$ 方向的力。$+\vect e_j$ 面的面力是 $\sigma_{ij}(\vect x+h\vect e_j)$，面积 $h^2$；$-\vect e_j$ 面的面力是 $-\sigma_{ij}(\vect x)$，面积也是 $h^2$。两块加起来，沿 $\vect e_i$ 的净力是
\[
\bigl[\sigma_{ij}(\vect x+h\vect e_j)-\sigma_{ij}(\vect x)\bigr]h^2
\;\approx\;h^3\,\partial_j\sigma_{ij}.
\]
除以体积 $h^3$，就得到\textbf{单位体积的净面力}的第 $i$ 个分量 $\partial_j\sigma_{ij}$。对 $j=1,2,3$ 求和，就是 $\nabla\cdot\ten\sigma$。所以“散度”不是又一个符号，而是“三个方向各查一次应力差，再相加”。

\textbf{一个能算的例子。}设
\[
\ten\sigma(x_1)=\frac{\sigma_0 x_1}{L}\,\vect e_1\otimes\vect e_1
\qquad(\sigma_0>0,\;L>0),
\]
即只有 $\sigma_{11}=\sigma_0x_1/L$，沿 $x_1$ 线性增大。则
\[
(\nabla\cdot\ten\sigma)_1=\partial_1\sigma_{11}=\frac{\sigma_0}{L},
\qquad
(\nabla\cdot\ten\sigma)_2=(\nabla\cdot\ten\sigma)_3=0 .
\]
也就是说，每单位体积有一个沿 $+\vect e_1$ 的净面力 $\sigma_0/L$：$+x_1$ 那一侧的面拉得比 $-x_1$ 那一侧更用力。要在静力学里平衡它，需要体力
\[
\vect f=-\nabla\cdot\ten\sigma=\Bigl(-\frac{\sigma_0}{L},0,0\Bigr),
\]
即一个沿 $-\vect e_1$ 的均匀体力。反过来，把例子\textbf{换成另一个应力场}：若整个应力分量处处相同，即 $\sigma_{11}(\vect x)\equiv\sigma_{\mathrm c}$（$\sigma_{\mathrm c}$ 是常数，不再是上面那个随 $x_1$ 线性增大的场），则 $\partial_1\sigma_{11}=0$，内部没有净面力，只有端面上的外力——这就是普通的均匀单轴拉伸。

\textbf{平衡方程。}把面力沿闭合曲面加起来，用高斯定理换成体积分。对张量场 $\ten\sigma$，分量形式是
\[
\oint_S\sigma_{ij}n_j\dd S=\int_V\partial_j\sigma_{ij}\dd V,
\qquad
(\nabla\cdot\ten\sigma)_i=\partial_j\sigma_{ij},
\]
其中 $\partial_j=\partial/\partial x_j$。与体力 $\vect f$ 平衡得（体力是作用在体积上的力，如重力、惯性力，$\vect f$ 是单位体积的力）
\[
\nabla\cdot\ten\sigma+\vect f=\vect0\quad\text{（静力学）},
\qquad
\nabla\cdot\ten\sigma+\vect f=\rho\vect a\quad\text{（动力学）},
\]
其中 $\rho$ 是质量密度、$\vect a$ 是加速度。写得开就是三个方向的方程：
\begin{align*}
\partial_1\sigma_{11}+\partial_2\sigma_{12}+\partial_3\sigma_{13}+f_1&=0,\\
\partial_1\sigma_{21}+\partial_2\sigma_{22}+\partial_3\sigma_{23}+f_2&=0,\\
\partial_1\sigma_{31}+\partial_2\sigma_{32}+\partial_3\sigma_{33}+f_3&=0.
\end{align*}
每个方程只含一个方向的力分量，所以叫“三个方向的力平衡”。

\warnbox{\textbf{指标顺序要盯住。}按 §6.1 的约定，$\sigma_{ij}$ 的第一个指标是\textbf{力的方向}，第二个是\textbf{面的法向}。面力是 $t_i=\sigma_{ij}n_j$，所以高斯定理缩并的是\textbf{第二个指标}：$\partial_j\sigma_{ij}$。若 $\ten\sigma$ 对称，$\partial_j\sigma_{ji}$ 与它逐分量相等；若 $\ten\sigma$ 不对称（例如有体力矩的介质），两者一般\textbf{不同}，而且不是简单差一个整体符号——它们是两个不同的缩并，不能互相替换。$\nabla\cdot\ten\sigma$ 是\textbf{矢量}（缩并一次），$\nabla\ten\sigma$ 才是三阶张量；两者差一个点，差一个阶。}

\warnbox{\textbf{选读（可跳过）：高斯定理为什么成立。}
这里只给一句话的图像，不影响后面；跳过后直接进入 §7。
把体积切成许多小盒子。每个小盒子有六个面，相邻盒子共享的面上的面力大小相等、方向相反，求和时成对抵消；最后只剩下最外层的面。另一方面，每个小盒子的“流出减流入”近似等于 $\partial_j\sigma_{ij}$ 乘体积。把小盒子加起来，就得到 $\oint_S\sigma_{ij}n_j\dd S=\int_V\partial_j\sigma_{ij}\dd V$。一维类比是 $\int_a^b f'(x)\dd x=f(b)-f(a)$：内部点成对抵消，只剩端点。严格证明只需把“近似”用中值定理写成可控制的小量，本稿不展开。}

\hdr{§7}{独立练习}

下面十二道题覆盖本稿的全部操作：并矢运算与不变量（P1）、惯量张量（P2、P4）、欧拉方程（P3、P12）、并矢的左右作用与张量判定（P5、P6）、迹与对称化（P7、P8）、惯量的构造（P9）、应力（P10）、旋转算子（P11）。
\textbf{用法：}先自己在纸上做，不看附录；做完再翻附录 A 对详解。正文只给题干与考点，不给关键步骤或结果。P4 用到 P2 的 $\ten J$，若还没做 P2，请先做 P2（或按 P2 的公式自己算一遍）。

\practice{P1}{自编题，承旧稿例 1}
{设 $\ten D=4\vect e_1\otimes\vect e_1+\vect e_1\otimes\vect e_2+\vect e_2\otimes\vect e_1+4\vect e_2\otimes\vect e_2+2\vect e_3\otimes\vect e_3$，$\vect x=\vect e_1+2\vect e_2+3\vect e_3$。求 $\ten D\cdot\vect x$、$\tr\ten D$、$\det\ten D$、$\ten D:\ten D$、$\ten D$ 的特征值与一组主轴，并写出谱分解。}
{并矢的分量、右作用、迹/行列式/双点积，以及对称张量的特征值与主轴。}{ans:P1}

\practice{P2}{自编题，承旧稿例 2}
{均匀长方体质量 $m=5$ kg，边长 $a=2$、$b=3$、$c=4$ m，质心在原点，棱沿坐标轴。求它的转动惯量张量 $\ten J$；再把它压成薄板（$c\to0$，质量仍为 $m$），验证垂直轴定理。}
{长方体惯量张量公式、薄板极限与垂直轴定理。}{ans:P2}

\practice{P3}{自编题，承旧稿例 3}
{对称陀螺 $A=B=2$ kg·m$^2$、$C=3$ kg·m$^2$，自转 $\omega_3=10$ rad/s，无外力矩。求横向分量绕体轴的角速度与周期，并说明转向。}
{欧拉方程在 $A=B$ 时的解、进动角速度与周期。}{ans:P3}

\practice{P4}{自编题，承旧稿例 4}
{用 P2 得到的 $\ten J$。求沿与三条棱成等角的方向 $\vect n=(1,1,1)/\sqrt3$ 的转动惯量 $I(\vect n)=\vect n\cdot\ten J\cdot\vect n$（这个方向不是 P2 长方体的空间对角线）；再取 $\boldsymbol\omega=\vect n$（$|\boldsymbol\omega|=1$），验证 $\boldsymbol\omega\cdot\vect L=2K$。}
{二次型沿方向的值、角动量与转动动能的关系。}{ans:P4}

\practice{P5}{自编概念题，承旧稿附录 B 第 1 问}
{判断并说明：(a) $(\vect a\otimes\vect b)\cdot\vect c$ 与 $\vect a(\vect b\cdot\vect c)$ 是否相同？(b) $(\vect a\otimes\vect b)\cdot\vect c$ 与 $(\vect a\cdot\vect b)\vect c$ 是否相同？}
{并矢右作用与两种“先点乘再乘矢量”的区别。}{ans:P5}

\practice{P6}{自编概念题，承旧稿附录 B 第 2 问}
{给出一组九个排成方阵的数，怎样判断它是不是二阶张量？举一个“在旋转下像矢量、但坐标系允许平移时就不像矢量”的例子。}
{张量的变换律与“对象/账本”的区别。}{ans:P6}

\practice{P7}{自编概念题，承旧稿附录 B 第 3 问}
{用 $\tr(\ten A\ten B)=\tr(\ten B\ten A)$ 写出两步，证明 $\tr(\ten R^{\mathsf T}\ten T\ten R)=\tr\ten T$，并说明每一步用了什么性质。}
{迹的循环性与正交换基下的不变性。}{ans:P7}

\practice{P8}{自编概念题，承旧稿附录 B 第 4 问}
{二次型 $\vect x\cdot\ten D\cdot\vect x$ 为什么只与 $\ten D$ 的对称部分有关？举一个 $2\times2$ 的数字例子验证。}
{指标重命名、对称化与二次型。}{ans:P8}

\practice{P9}{自编概念题，承旧稿附录 B 第 5 问}
{$J_{12}=-\int xy\dd m$ 里的负号从哪里来？如果把它丢掉，矩阵还对称吗？谱定理还成立吗？算出的主转动惯量还是真实的吗？}
{惯量张量的构造、对称性与主值的物理意义。}{ans:P9}

\practice{P10}{自编概念题，承旧稿附录 B 第 6 问}
{单轴拉伸 $\ten\sigma=\sigma_0\vect e_1\otimes\vect e_1$。最大剪应力多大？出现在哪个截面？为什么这个结果不需要背？}
{面力 $\vect t=\ten\sigma\cdot\vect n$、斜截面分解与最大剪应力。}{ans:P10}

\practice{P11}{自编题}
{设 $S(t)$ 是绕 $\vect e_3$ 轴转 $\varphi(t)=\alpha t$ 的旋转算子（$\alpha$ 为常数）。(a) 写出 $S$ 在 $\{\vect e_1,\vect e_2,\vect e_3\}$ 下的矩阵；(b) 验证 $S^{\mathsf T}S=\unitd$；(c) 用 $\boldsymbol\Omega^{\times}=(\dd S/\dd t)S^{\mathsf T}$ 求 $\boldsymbol\Omega$。}
{旋转算子的矩阵、正交性与角速度。}{ans:P11}

\practice{P12}{自编题}
{刚体的三个主转动惯量为 $A=1$、$B=2$、$C=3$ kg·m$^2$。某瞬时在体坐标系中 $\boldsymbol\omega=(1,1,1)$ rad/s，外力矩 $\vect M=\vect0$。用欧拉方程求此刻 $\dot\omega_1,\dot\omega_2,\dot\omega_3$。}
{欧拉方程的分量形式与陀螺项。}{ans:P12}

\hdr{§8}{易错点、速查与三句话}

\warnbox{\textbf{八个高频错误}\\[2pt]
1. \textbf{把张量当成“九个数字”}：九个数字要按变换律 $T'_{ij}=R_{ki}R_{lj}T_{kl}$ 变才叫张量（§2.7）。不检验就套用谱定理，是错的源头。\\
2. \textbf{把 $\vect a\otimes\vect b$ 当成 $\vect a\cdot\vect b$}：前者是算子（九个分量），后者是标量。中间是 $\otimes$ 还是 $\cdot$，差一个层级。\\
3. \textbf{左右的顺序混用}：作为算子，$\ten D\cdot\vect x$ 与 $\vect x\cdot\ten D$ 一般不同；它们对\textbf{所有} $\vect x$ 相同，当且仅当 $\ten D$ 对称（§2.3）。力学里 $\ten J$、$\ten\sigma$ 对称，所以“碰巧”怎么写都对；换到一般旋转算子 $S$ 上，左右一混通常出错（对某个特定输入仍可能碰巧相同）。\\
4. \textbf{直接写 $\dd\vect L/\dd t=\ten J\cdot\dot{\boldsymbol\omega}$}：只在 $\ten J$ 于该系为常量时成立；体坐标系里求的是体导数，固定系里还要补 $\boldsymbol\omega\times\vect L$（§5.1）。\\
5. \textbf{把惯量张量的非对角元符号写反}：$J_{12}=-\int xy\dd m$，负号不能丢；丢了以后主轴会算错（P9）。\\
6. \textbf{把 $\nabla\cdot\ten\sigma$ 的指标缩并错}：按“力方向在前、法向在后”的约定，面力是 $t_i=\sigma_{ij}n_j$，所以散度是 $\partial_j\sigma_{ij}$；对称时与 $\partial_j\sigma_{ji}$ 相同，非对称时不同（§6.4）。\\
7. \textbf{把特征值当成“唯一的解”}：三个特征值不变，但\textbf{特征向量的分量随基变}（§1.6）。说“主轴方向”必须连着坐标系或刚体说。\\
8. \textbf{把换基矩阵的转置写反}：$R$ 的列是新基在旧基下的分量时，矢量的新坐标是 $\vect x'=\ten R^{\mathsf T}\vect x$，二阶张量是 $\ten T'=\ten R^{\mathsf T}\ten T\ten R$（§0.4）；不是 $\ten R\vect x$。\\
}

\begin{center}
\tabtop
\begin{tabular}{@{}p{0.32\linewidth}p{0.58\linewidth}@{}}
\thead{要做的事} & \thead{用什么} \\
\multicolumn{2}{@{}l@{}}{\color{petal}\rule{\linewidth}{0.5pt}}\\
判断是不是张量 & 查变换律 $T'_{ij}=R_{ki}R_{lj}T_{kl}$；或对已知矢量的线性关系用商法则（§2.7）\\[3pt]
把“矢量 $\to$ 矢量”写成式子 & $\ten D=D_{ij}\vect e_i\otimes\vect e_j$，作用规则 $(\vect a\otimes\vect b)\cdot\vect x=\vect a(\vect b\cdot\vect x)$\\[3pt]
判断是否对称 & 看 $D_{ij}=D_{ji}$；反对称部分等价于一个轴矢量（叉乘，§2.6）\\[3pt]
求主轴与主值 & 解 $\det(\ten D-\lambda\unitd)=0$；对称时主轴正交，$\ten D=\sum_k\lambda_k\vect n_k\otimes\vect n_k$（§2.9）\\[3pt]
验算 & $I_1=\tr$、$I_2$、$I_3=\det$；惯量张量还有 $A+B\ge C$（§1.6、§4.3）\\[3pt]
转动惯量 & $\ten J=\int(r^2\unitd-\vect r\otimes\vect r)\dd m$；平行轴 $\ten J_O=\ten J_C+m(R^2\unitd-\vect R\otimes\vect R)$（§4.1、§4.4）\\[3pt]
角动量与动能 & $\vect L=\ten J\cdot\boldsymbol\omega$，$K=\tfrac12\boldsymbol\omega\cdot\ten J\cdot\boldsymbol\omega$，$\boldsymbol\omega\cdot\vect L=2K$（§4.2）\\[3pt]
体坐标系动力学 & $(\dd\vect L/\dd t)_{\text{体}}+\boldsymbol\omega\times\vect L=\vect M$，即欧拉方程（§5.2）\\[3pt]
面力 & $\vect t(\vect n)=\ten\sigma\cdot\vect n$；$\sigma_{ij}=\sigma_{ji}$；主应力 = 特征值（§6.1--§6.2）\\
\end{tabular}
\tabbot
\end{center}

\flowerrule

\subhdr{三句话总结}
\begin{enumerate}[label=\arabic*.,leftmargin=1.6em]
\item \textbf{张量是对象，矩阵是账本。}换基时账本变、对象不变；迹、行列式、特征值是账本里不随基变的数。
\item \textbf{并矢只是二阶张量的另一种写法。}作用规则 $(\vect a\otimes\vect b)\cdot\vect x=\vect a(\vect b\cdot\vect x)$ 一条就够；左右作用作为算子一般不对称，只有张量对称时才对所有输入相同；双点积吃掉两个指标、谱分解是三个投影并矢相加。
\item \textbf{理论力学里三处落点用的是同一台机器。}惯量张量给出 $\vect L=\ten J\cdot\boldsymbol\omega$，欧拉方程是它在体坐标系里的求导，应力张量是“面力对法向线性”的商法则产物。
\end{enumerate}

\clearpage
\appendix
\hdr{附录 A}{练习详解}

\answerstart{P1}{自编题，承旧稿例 1}
{设 $\ten D=4\vect e_1\otimes\vect e_1+\vect e_1\otimes\vect e_2+\vect e_2\otimes\vect e_1+4\vect e_2\otimes\vect e_2+2\vect e_3\otimes\vect e_3$，$\vect x=\vect e_1+2\vect e_2+3\vect e_3$。求 $\ten D\vect x$、$\tr\ten D$、$\det\ten D$、$\ten D:\ten D$、特征值与主轴，并写出谱分解。}

先把九个系数按行列摆好（第 $i$ 行第 $j$ 列是 $\vect e_i\otimes\vect e_j$ 的系数）：
\[
\ten D=\begin{pmatrix}4&1&0\\1&4&0\\0&0&2\end{pmatrix}.
\]
作用：$(\ten D\vect x)_i=\sum_jD_{ij}x_j$。第一行 $4\cdot1+1\cdot2+0=6$，第二行 $1\cdot1+4\cdot2=9$，第三行 $0+0+2\cdot3=6$，所以
\[
\ten D\vect x=(6,9,6).
\]
因为 $\ten D$ 对称，左作用 $\vect x\cdot\ten D$ 给同一个结果（§2.3）。

迹与行列式：$\tr\ten D=4+4+2=10$；按第三行展开，$\det\ten D=2(16-1)=30$。
双点积：$\ten D:\ten D=\sum_{i,j}D_{ij}^2=16+1+1+16+4=38$。

特征方程由 §0.5：$-\lambda^3+10\lambda^2-I_2\lambda+30=0$，其中 $I_2=\tfrac12(10^2-38)=31$。即
\[
\lambda^3-10\lambda^2+31\lambda-30=0 .
\]
试根 $\lambda=2$：$8-40+62-30=0$，成立。分解
\[
(\lambda-2)(\lambda^2-8\lambda+15)=(\lambda-2)(\lambda-3)(\lambda-5)=0,
\]
所以特征值是 $2,3,5$。

主轴：$\lambda=5$ 时
\[
(\ten D-5\unitd)=\begin{pmatrix}-1&1&0\\1&-1&0\\0&0&-3\end{pmatrix},
\]
前两行都给出 $-x+y=0$，即 $x=y$；第三行给 $-3z=0$，即 $z=0$。所以方向是 $(1,1,0)$，单位化 $\vect n_1=(\vect e_1+\vect e_2)/\sqrt2$。代回核验：
\[
(\ten D-5\unitd)\begin{pmatrix}1\\1\\0\end{pmatrix}
=\begin{pmatrix}-1+1\\1-1\\0\end{pmatrix}=\vect0 .
\]
$\lambda=3$ 时 $(\ten D-3\unitd)$ 的前两行给 $x+y=0$，取 $\vect n_2=(\vect e_1-\vect e_2)/\sqrt2$（也可以取相反方向，差一个符号）；$\lambda=2$ 时前两行给 $2x+y=0$ 与 $x+2y=0$，只能 $x=y=0$，第三行不限制 $z$，所以 $\vect n_3=\vect e_3$。谱分解
\[
\ten D=5\,\vect n_1\otimes\vect n_1+3\,\vect n_2\otimes\vect n_2+2\,\vect e_3\otimes\vect e_3 .
\]
验算（§1.6）：$5+3+2=10=\tr$，$5\cdot3\cdot2=30=\det$，$25+9+4=38=\ten D:\ten D$，都对。

\answerstart{P2}{自编题，承旧稿例 2}
{均匀长方体 $m=5$ kg，$a=2,b=3,c=4$ m，质心在原点，棱沿坐标轴。求 $\ten J$；再压成薄板（$c\to0$，质量仍为 $m$），验证垂直轴定理。}

由 §4.4 直接代入：
\[
\begin{aligned}
\ten J&=\frac{m}{12}\operatorname{diag}(b^2+c^2,\;a^2+c^2,\;a^2+b^2)\\
&=\frac{5}{12}\operatorname{diag}(25,\;20,\;13)
=\operatorname{diag}\Bigl(\frac{125}{12},\frac{100}{12},\frac{65}{12}\Bigr)\\
&\approx\operatorname{diag}(10.417,\;8.333,\;5.417)\ \text{kg·m}^2 .
\end{aligned}
\]
验算三角形不等式：$10.417+8.333=18.750\ge5.417$；$10.417+5.417=15.834\ge8.333$；$8.333+5.417=13.750\ge10.417$，三个循环都成立（§4.3）。

\answercontinue{P2}{P2 的 $\ten J$ 已算出；本页接续做薄板极限，并验证 $J_z=J_x+J_y$。}
压成薄板（$c=0$，质量仍为 $5$ kg）后
\[
\ten J_{\text{板}}=\frac{m}{12}\operatorname{diag}(b^2,\;a^2,\;a^2+b^2)
=\frac{5}{12}\operatorname{diag}(9,\;4,\;13)
=\operatorname{diag}(3.750,\;1.667,\;5.417).
\]
于是 $J_z=65/12$，而 $J_x+J_y=(45+20)/12=65/12$，二者相等，垂直轴定理成立。

\answerstart{P3}{自编题，承旧稿例 3}
{对称陀螺 $A=B=2$ kg·m$^2$、$C=3$ kg·m$^2$，自转 $\omega_3=10$ rad/s，无外力矩。求横向分量绕体轴的角速度与周期，并说明转向。}

由 §5.3，
\[
\Omega=\frac{C-A}{A}\omega_3=\frac{3-2}{2}\times10=5\ \text{rad/s},
\qquad
T=\frac{2\pi}{\Omega}=\frac{2\pi}{5}\approx1.257\ \text{s}.
\]
即 $\boldsymbol\omega$ 的横向分量每 $1.257$ 秒绕体轴转一圈，$\omega_3$ 与 $|\boldsymbol\omega|$ 都不变。因为 $C>A$，$\Omega>0$，所以横向分量沿 $+\vect e_3$ 的正向（右手定则）转；在体坐标系里看，$\boldsymbol\omega$ 的端点绕 $\vect e_3$ 画圆，$\vect L$ 与 $\boldsymbol\omega$、$\vect e_3$ 始终共面（§5.3 的图 7）。

\answerstart{P4}{自编题，承旧稿例 4}
{用 P2 的 $\ten J$。求沿与三条棱成等角的方向 $\vect n=(1,1,1)/\sqrt3$ 的转动惯量 $I(\vect n)=\vect n\cdot\ten J\cdot\vect n$（该方向不是 P2 长方体的空间对角线 $(2,3,4)/\sqrt{29}$）；再取 $\boldsymbol\omega=\vect n$，验证 $\boldsymbol\omega\cdot\vect L=2K$。}

P2 已得
\[
\ten J=\operatorname{diag}\Bigl(\frac{125}{12},\frac{100}{12},\frac{65}{12}\Bigr).
\]
沿单位方向 $\vect n$ 的转动惯量是二次型在 $\vect n$ 上的值（§1.5 的瑞利商）：
\[
I(\vect n)=\vect n\cdot\ten J\vect n
=\frac{1}{3}\Bigl(\frac{125}{12}+\frac{100}{12}+\frac{65}{12}\Bigr)
=\frac{1}{3}\cdot\frac{290}{12}
=\frac{145}{18}\approx8.056\ \text{kg·m}^2 .
\]
它落在最小主值 $65/12\approx5.417$ 与最大主值 $125/12\approx10.417$ 之间，符合 $\lambda_{\min}\le I\le\lambda_{\max}$。

\answercontinue{P4}{P4 已得 $I(\vect n)=145/18$（$\ten J=\operatorname{diag}(125/12,100/12,65/12)$，$\vect n=(1,1,1)/\sqrt3$）；本页验证 $\boldsymbol\omega=\vect n$ 时的 $\boldsymbol\omega\cdot\vect L=2K$。}

验证：取 $\boldsymbol\omega=\vect n$（$|\boldsymbol\omega|=1$），则 $K=\tfrac12\boldsymbol\omega\cdot\ten J\boldsymbol\omega=\tfrac12I=145/36$；
而
\[
\vect L=\ten J\boldsymbol\omega=\frac{1}{\sqrt3}\Bigl(\frac{125}{12},\frac{100}{12},\frac{65}{12}\Bigr),
\qquad
\boldsymbol\omega\cdot\vect L=\frac{1}{3}\cdot\frac{290}{12}=\frac{145}{18}=2K .
\]
最后这个等式是恒等式（不依赖是否在主轴系）：$\boldsymbol\omega\cdot(\ten J\boldsymbol\omega)=\boldsymbol\omega\cdot\ten J\cdot\boldsymbol\omega=2K$。

\answerstart{P5}{自编概念题，承旧稿附录 B 第 1 问}
{判断并说明：(a) $(\vect a\otimes\vect b)\cdot\vect c$ 与 $\vect a(\vect b\cdot\vect c)$ 是否相同？(b) $(\vect a\otimes\vect b)\cdot\vect c$ 与 $(\vect a\cdot\vect b)\vect c$ 是否相同？}

(a) 相同。按右作用规则 $(\vect a\otimes\vect b)\cdot\vect c=\vect a(\vect b\cdot\vect c)$；读法是“先把 $\vect c$ 投影到 $\vect b$ 上得到数 $\vect b\cdot\vect c$，再沿 $\vect a$ 放大”。

(b) 一般不同。左边是“投到 $\vect b$、沿 $\vect a$ 放大”；右边是“先算 $\vect a\cdot\vect b$ 这个数，再沿 $\vect c$ 放大”。
取 $\vect a=\vect e_1$、$\vect b=\vect e_1$、$\vect c=\vect e_2$：
\[
(\vect e_1\otimes\vect e_1)\cdot\vect e_2=\vect e_1(\vect e_1\cdot\vect e_2)=\vect0,
\qquad
(\vect e_1\cdot\vect e_1)\vect e_2=\vect e_2 .
\]
一个为零，一个不为零，所以不相同。两者只有在特殊输入下才碰巧相等。例如让 $\vect c$ 沿 $\vect a$，即 $\vect c=k\vect a$：
\[
(\vect a\otimes\vect b)\cdot(k\vect a)=k\vect a(\vect b\cdot\vect a),
\qquad
(\vect a\cdot\vect b)(k\vect a)=k\vect a(\vect a\cdot\vect b),
\]
两边相同（数乘可以交换）。这说明“一般不同、特殊输入可以相同”，但不能因此把两个表达式当成同一个。

\answerstart{P6}{自编概念题，承旧稿附录 B 第 2 问}
{给出一组九个排成方阵的数，怎样判断它是不是二阶张量？举一个“在旋转下像矢量、但坐标系允许平移时就不像矢量”的例子。}

判据是\textbf{变换律}：对任意正交换基（$R_{ki}=\vect n_i\cdot\vect e_k$ 是新的正交基），检查
\[
T'_{ij}=R_{ki}R_{lj}T_{kl}
\]
是否成立。成立才是二阶张量；不成立只是九个数字。实际操作中，若已知一个物理关系 $\vect u=\ten T\vect v$ 对任意矢量 $\vect v$ 成立，且 $\vect u,\vect v$ 都是矢量，则可用商法则（§2.7）直接判定 $\ten T$ 是张量，不必逐项验证。

例子：$\delta_{ij}$ 与 $a_ib_j$ 都是张量，前者已在 §2.7 验过，后者就是并矢（§2.2）。
\textbf{位置矢量} $(x_1,x_2,x_3)$ 在旋转下按 $\vect x'=\ten R^{\mathsf T}\vect x$ 变，所以对旋转来说它是矢量；但若允许坐标系平移 $\vect x\to\vect x+\vect a$，分量会多出一个常数项，不再是齐次的，所以对“平移 + 旋转”这个更大的群它不是矢量。\textbf{“是不是张量”必须连着允许哪些变换说}；本稿只允许旋转与反射，在这个群里位置矢量是矢量。

\answerstart{P7}{自编概念题，承旧稿附录 B 第 3 问}
{用 $\tr(\ten A\ten B)=\tr(\ten B\ten A)$ 写出两步，证明 $\tr(\ten R^{\mathsf T}\ten T\ten R)=\tr\ten T$，并说明每一步用了什么性质。}

两步：
\[
\tr(\ten R^{\mathsf T}\ten T\ten R)
=\tr\bigl(\ten R\ten R^{\mathsf T}\ten T\bigr)
=\tr\ten T .
\]
第一步把 $\ten R$ 从最右边“转圈”到最左边，用的是迹的循环性 $\tr(\ten A\ten B)=\tr(\ten B\ten A)$（把 $\ten A=\ten R^{\mathsf T}\ten T$、$\ten B=\ten R$ 代入）；第二步用正交性 $\ten R\ten R^{\mathsf T}=\unitd$。如果愿意写细一点，也可以写成
\[
\tr\bigl((\ten R^{\mathsf T}\ten T)\ten R\bigr)
=\tr\bigl(\ten R(\ten R^{\mathsf T}\ten T)\bigr)
=\tr\bigl((\ten R\ten R^{\mathsf T})\ten T\bigr)
=\tr\ten T,
\]
中间两步分别是循环性与结合律。两种写法是同一件事。

\answerstart{P8}{自编概念题，承旧稿附录 B 第 4 问}
{二次型 $\vect x\cdot\ten D\cdot\vect x$ 为什么只与 $\ten D$ 的对称部分有关？举一个 $2\times2$ 的数字例子验证。}

把二次型写成 $\sum_{i,j}D_{ij}x_ix_j$。交换 $i,j$ 的名字（只是换字母，值不变）：
\[
\sum_{i,j}D_{ij}x_ix_j=\sum_{i,j}D_{ji}x_jx_i=\sum_{i,j}D_{ji}x_ix_j .
\]
两式相加除二，得到 $\sum_{i,j}\tfrac12(D_{ij}+D_{ji})x_ix_j$，只含对称部分 $S_{ij}=\tfrac12(D_{ij}+D_{ji})$。所以反对称部分对二次型没有贡献。

数字例子：取 $\ten D=\begin{pmatrix}0&3\\0&0\end{pmatrix}$，则
\[
E=\begin{pmatrix}x&y\end{pmatrix}\begin{pmatrix}0&3\\0&0\end{pmatrix}\begin{pmatrix}x\\y\end{pmatrix}=3xy .
\]
\answercontinue{P8}{P8 的数字例子已算到 $E=3xy$；本页接续算对称部分，并比较两个二次型。}
对称部分是 $\ten S=\begin{pmatrix}0&1.5\\1.5&0\end{pmatrix}$，它给出的二次型是
\[
E=\begin{pmatrix}x&y\end{pmatrix}\begin{pmatrix}0&1.5\\1.5&0\end{pmatrix}\begin{pmatrix}x\\y\end{pmatrix}=1.5xy+1.5yx=3xy,
\]
一样。注意原矩阵不对称，但二次型只“看见”它的对称部分。

\answerstart{P9}{自编概念题，承旧稿附录 B 第 5 问}
{$J_{12}=-\int xy\dd m$ 里的负号从哪里来？如果把它丢掉，矩阵还对称吗？谱定理还成立吗？算出的主转动惯量还是真实的吗？}

负号来自 §4.1 的构造：$\ten J=\int(r^2\unitd-\vect r\otimes\vect r)\dd m$，而 $\vect r\otimes\vect r$ 的分量是 $x_ix_j$。对 $i=1,j=2$，$r^2\delta_{12}=0$，所以
\[
J_{12}=\int(0-x_1x_2)\dd m=-\int xy\dd m .
\]
它不是“人为加的”，是 $-\vect r\otimes\vect r$ 那一项的直接结果。物理上它保证“沿 $\vect r$ 方向转时惯量为零”这类事实（P9 下面的例子）。

如果只把 $J_{12}=J_{21}$ 的负号去掉（即把它们写成 $+\int xy\dd m$，其余分量不动），矩阵仍然是对称的，§1.5 的谱定理也仍然成立（对称矩阵总能正交对角化），但主转动惯量会变。举一个让非对角元真正起作用的例子：一个质点 $m$ 在 $\vect r=(1,1,1)$。真实的
\[
\ten J=m(r^2\unitd-\vect r\otimes\vect r)
=m\begin{pmatrix}2&-1&-1\\-1&2&-1\\-1&-1&2\end{pmatrix},
\]
它的三个主值是 $0,3m,3m$：沿 $\vect r$ 方向（质点就在这条直线上）惯量为零，垂直于 $\vect r$ 的任意方向为 $3m$。若把 $J_{12}=J_{21}=-m$ 错写成 $+m$，得到
\[
\ten J'=m\begin{pmatrix}2&1&-1\\1&2&-1\\-1&-1&2\end{pmatrix},
\]
它仍对称、仍可正交对角化。注意 $\ten J'$ 带整体因子 $m$：把它写成 $\ten J'=m\ten A$，其中
\[
\ten A=\begin{pmatrix}2&1&-1\\1&2&-1\\-1&-1&2\end{pmatrix}.
\]
整体因子只把特征值乘上 $m$、不改变特征向量，所以 $\ten A$ 的特征方程是 $(\lambda-1)^2(\lambda-4)=0$，$\ten J'$ 的主值就是 $m,m,4m$（等价地，直接解 $\det(\ten J'-\lambda\unitd)=0$ 也会得到 $\lambda=m,m,4m$）。对比 $0,3m,3m$，可见“对称”只保证谱定理可用，不保证主值正确。负号是构造式 $-\vect r\otimes\vect r$ 的一部分，丢了它就换了另一个张量。（另一种更粗的错误是把整个构造写成 $r^2\unitd+\vect r\otimes\vect r$；那会连对角元一起改，后果更大。）\textbf{所以“对称”不等于“正确”；物理内容要靠构造式保证。}

\answerstart{P10}{自编概念题，承旧稿附录 B 第 6 问}
{单轴拉伸 $\ten\sigma=\sigma_0\vect e_1\otimes\vect e_1$。最大剪应力多大？出现在哪个截面？为什么这个结果不需要背？}

取法向与 $\vect e_1$ 成 $\theta$ 的截面（在 $\vect e_1\vect e_2$ 面内），$\vect n=\cos\theta\,\vect e_1+\sin\theta\,\vect e_2$。面力
\[
\vect t=\ten\sigma\cdot\vect n=\sigma_0\cos\theta\,\vect e_1 .
\]
法向分量 $t_n=\vect t\cdot\vect n=\sigma_0\cos^2\theta$；切向分量大小
\[
\tau=\sqrt{|\vect t|^2-t_n^2}
=\sigma_0\sin\theta\cos\theta
=\frac{\sigma_0}{2}\sin2\theta .
\]
所以 $\tau$ 在 $\sin2\theta=1$ 时最大，即 $2\theta=90^\circ$、$\theta=45^\circ$，最大值
\[
\tau_{\max}=\frac{\sigma_0}{2}.
\]
\answercontinue{P10}{P10 已算出 $\tau=\sigma_0\sin\theta\cos\theta=\tfrac{\sigma_0}{2}\sin2\theta$；本页接续说明最大值出现在哪个截面。}
出现在法向与 $\vect e_1$ 成 $45^\circ$ 的截面上（例如 $\vect n=(\vect e_1+\vect e_2)/\sqrt2$ 或 $(\vect e_1-\vect e_2)/\sqrt2$；绕 $\vect e_1$ 转一圈的任一这样的截面都一样）。在 $45^\circ$ 截面上，法向应力也正好是 $\sigma_0/2$。

\textbf{为什么不需要背：}只要写出 $\vect t=\ten\sigma\cdot\vect n$，再把它分成法向与切向，$\tau=\sigma_0\sin\theta\cos\theta$ 就出来了；最大值是 $\sin2\theta$ 的最大值，不需要记特殊角度。这个例子也解释了“拉伸试件常沿 $45^\circ$ 断裂”：那个截面上剪应力最大。

\answerstart{P11}{自编题}
{设 $S(t)$ 是绕 $\vect e_3$ 轴转 $\varphi(t)=\alpha t$ 的旋转算子（$\alpha$ 为常数）。(a) 写出 $S$ 的矩阵；(b) 验证 $S^{\mathsf T}S=\unitd$；(c) 用 $\boldsymbol\Omega^{\times}=(\dd S/\dd t)S^{\mathsf T}$ 求 $\boldsymbol\Omega$。}

(a) 由 §3.2，绕 $\vect e_3$ 转 $\varphi$ 的矩阵是
\[
S(t)=
\begin{pmatrix}
\cos\alpha t&-\sin\alpha t&0\\
\sin\alpha t&\cos\alpha t&0\\
0&0&1
\end{pmatrix}.
\]
它的第 $1$ 列是 $\vect n'_1=(\cos\alpha t,\sin\alpha t,0)$，第 $2$ 列是 $\vect n'_2=(-\sin\alpha t,\cos\alpha t,0)$，第 $3$ 列是 $\vect e_3$。

(b) 记 $c=\cos\alpha t$、$s=\sin\alpha t$。直接乘：
\[
S^{\mathsf T}S=
\begin{pmatrix}
c&s&0\\-s&c&0\\0&0&1
\end{pmatrix}
\begin{pmatrix}
c&-s&0\\s&c&0\\0&0&1
\end{pmatrix}
=
\begin{pmatrix}
c^2+s^2&-cs+sc&0\\
-sc+cs&s^2+c^2&0\\
0&0&1
\end{pmatrix}
=\unitd .
\]
同理 $SS^{\mathsf T}=\unitd$，所以 $S$ 是正交算子（§3.3）。

\answercontinue{P11}{题设仍是绕 $\vect e_3$ 转 $\varphi=\alpha t$；上页已确认 $S^{\mathsf T}S=SS^{\mathsf T}=\unitd$；本页做 (c) 求导，记 $c=\cos\alpha t$、$s=\sin\alpha t$。}
(c) 求导：
\[
\frac{\dd S}{\dd t}=\alpha
\begin{pmatrix}
-\sin\alpha t&-\cos\alpha t&0\\
\cos\alpha t&-\sin\alpha t&0\\
0&0&0
\end{pmatrix}.
\]
于是
\[
\frac{\dd S}{\dd t}S^{\mathsf T}
=\alpha
\begin{pmatrix}
-s&-c&0\\ c&-s&0\\0&0&0
\end{pmatrix}
\begin{pmatrix}
c&s&0\\-s&c&0\\0&0&1
\end{pmatrix}
=\alpha
\begin{pmatrix}
0&-1&0\\1&0&0\\0&0&0
\end{pmatrix}
=\alpha\,\vect e_3^{\times}.
\]
最后一个等号因为
\[
\vect e_3^{\times}=
\begin{pmatrix}0&-1&0\\1&0&0\\0&0&0\end{pmatrix}.
\]
所以 $\boldsymbol\Omega^{\times}=\alpha\vect e_3^{\times}$，即
\[
\boldsymbol\Omega=\alpha\,\vect e_3 .
\]
这与“绕 $\vect e_3$ 以角速度 $\alpha$ 转”一致；方向沿转轴，大小是转角的变化率。

\answerstart{P12}{自编题}
{刚体 $A=1$、$B=2$、$C=3$ kg·m$^2$。某瞬时在体坐标系中 $\boldsymbol\omega=(1,1,1)$ rad/s，外力矩 $\vect M=\vect0$。用欧拉方程求此刻 $\dot\omega_1,\dot\omega_2,\dot\omega_3$。}

由 §5.2 的欧拉方程（$M_k=0$）：
\[
\begin{aligned}
A\dot\omega_1+(C-B)\omega_2\omega_3&=0,\\
B\dot\omega_2+(A-C)\omega_3\omega_1&=0,\\
C\dot\omega_3+(B-A)\omega_1\omega_2&=0 .
\end{aligned}
\]
代入 $A=1,B=2,C=3$、$\omega_1=\omega_2=\omega_3=1$：
\begin{gather*}
1\cdot\dot\omega_1+(3-2)\cdot1\cdot1=0 \quad\Longrightarrow\quad \dot\omega_1=-1,\\
2\cdot\dot\omega_2+(1-3)\cdot1\cdot1=0 \quad\Longrightarrow\quad 2\dot\omega_2-2=0 \quad\Longrightarrow\quad \dot\omega_2=1,\\
3\cdot\dot\omega_3+(2-1)\cdot1\cdot1=0 \quad\Longrightarrow\quad 3\dot\omega_3+1=0 \quad\Longrightarrow\quad \dot\omega_3=-\frac13 .
\end{gather*}
所以
\[
\dot{\boldsymbol\omega}=\Bigl(-1,\;1,\;-\frac13\Bigr)\ \text{rad/s}^2 .
\]
核对：把三个结果代回三条方程，左边分别是 $1\cdot(-1)+1=0$、$2\cdot1-2=0$、$3\cdot(-1/3)+1=0$，都对。注意这是一个瞬时的变化率；下一时刻 $\boldsymbol\omega$ 变了，右边的陀螺项也会变。

\hdr{附录 B}{指标与记号速查}

\aside{求和约定：同一项里重复出现两次的指标自动对 $1,2,3$ 求和。本稿正文常写求和号，附录与推导里用约定，避免“缩并了哪个指标”看错。
$\varepsilon_{ijk}$ 的取值见 §2.6：偶排列 $+1$、奇排列 $-1$、有重复 $0$。RIT Whelan notes 10 的对偶张量定义为 $\ast A=-\vect a^{\times}$（符号相反），引用时先对矩阵。}

\begin{center}
\begin{minipage}{\linewidth}
\small
\renewcommand{\arraystretch}{1.05}
\tabtop
\begin{tabular}{@{}ll@{}}
\thead{对象} & \thead{指标写法} \\
\multicolumn{2}{@{}l@{}}{\color{petal}\rule{0.55\linewidth}{0.5pt}}\\
并矢 & $D_{ij}=a_ib_j$ \\
右作用 & $(\ten D\cdot\vect x)_i=D_{ij}x_j$ \\
左作用 & $(\vect x\cdot\ten D)_j=x_iD_{ij}$ \\
张量乘张量 & $(\ten A\cdot\ten B)_{ij}=A_{ik}B_{kj}$ \\
双点积 & $\ten A:\ten B=A_{ij}B_{ij}$ \\
迹 & $\tr\ten D=D_{ii}$ \\
二次型 & $E=a_iD_{ij}a_j$ \\
变换律 & $D'_{ij}=R_{ki}R_{lj}D_{kl}$ \\
叉乘 & $(\vect a\times\vect b)_i=\varepsilon_{ijk}a_jb_k$ \\
对偶张量 & $(\vect a^{\times})_{ij}=\varepsilon_{ikj}a_k$ \\
转动惯量 & $J_{ij}=\int(r^2\delta_{ij}-x_ix_j)\dd m$ \\
散度 & $(\nabla\cdot\ten D)_i=\partial_jD_{ij}$ \\
\end{tabular}
\tabbot
\end{minipage}
\end{center}

\hdr{附录 C}{来源对应、补充与选读清单}

\textbf{一句话：}本稿以课件 p8--p12 为主要来源；惯量、欧拉方程、应力三节是课件未展开的延伸，本稿用同一套并矢语言补足；旧稿（2026-11-05）只作为主题范围与题目来源，章节、例子与图均重新组织。

\subhdr{C.1\quad 课件页与本稿位置}

\begin{center}
\tabtop
\begin{tabular}{@{}p{0.16\linewidth}p{0.34\linewidth}p{0.40\linewidth}@{}}
\thead{课件页} & \thead{内容} & \thead{本稿位置} \\
\multicolumn{3}{@{}l@{}}{\color{petal}\rule{\linewidth}{0.5pt}}\\
p8 & 正交基、求和约定、点乘、叉乘、$\varepsilon_{ijk}$、对偶张量 & §0.1、§2.6、附录 B \\
p9 & 张量积（Dyad）、左右作用、转置、线性算子、二次型、单位张量 & §1.1、§2.2--§2.5、§2.9、§2.10 \\
p10 & 固定基与随体基、$S(t)=\sum_k\vect n'_k\otimes\vect e_k$ & §3.1、§3.2 \\
p11 & $S^{\mathsf T}S=SS^{\mathsf T}=\unitd$、绕轴转 $\varphi$ 的公式、无穷小转动 & §3.3、§3.4、§3.5 \\
p12 & $\boldsymbol\Omega^{\times}=(\dd S/\dd t)S^{\mathsf T}$、$\dd\vect n'_k/\dd t=\boldsymbol\Omega\times\vect n'_k$ & §3.5、§5.1 \\
p13--p18 & 欧拉角、运动学例题（E8.01、E8.02、HW8.01--HW8.03） & 不属于本专题的并矢主线，未重复；需要时按课件与运动学讲义处理 \\
\end{tabular}
\tabbot
\end{center}

\needspace{8\baselineskip}
\subhdr{C.2\quad 旧稿主题与本稿位置}

\begin{center}
\tabtop
\begin{tabular}{@{}p{0.30\linewidth}p{0.54\linewidth}@{}}
\thead{旧稿主题} & \thead{本稿位置与处理} \\
\multicolumn{2}{@{}l@{}}{\color{petal}\rule{\linewidth}{0.5pt}}\\
旧 §1--§2（矩阵、换基、迹、行列式、特征多项式） & 重组为 §0.2--§0.5；换基改为“三步推导 + 核对”，并明确 $R$ 的列约定 \\
旧 §3（二次型、椭球、特征值、谱定理） & 重组为 §1；$45^\circ$ 例的坐标与图轴标签统一重算；谱定理证明标为选读 \\
旧 §4（阶、张量积、指标、变换律、商法则、上下标） & §2.1、§2.2、§2.7；上下标移入选读（§2.10） \\
旧 §5（并矢、左右作用、点积与双点积、单位与投影、转置与对称分解、对偶张量、缩并表、谱分解） & §2.3--§2.9；缩并表作为查阅表；谱分解与 §1.4 的例子对上 \\
旧 §6（三记号字典） & §2.10，保留并补充“并矢不是矩阵”的说明 \\
旧 §7（惯量张量） & §4，补出正定性、三角形不等式、对称轴判断与平行轴推导；椭球例标为选读 \\
旧 §8（输运定理、欧拉方程、对称陀螺） & §5，补出并矢版输运定理的逐项推导与两个核对；对称陀螺例换成不同数字 \\
旧 §9（应力张量） & §6，补出柯西定理的四面体论证、指标顺序、Gauss 定理的说明 \\
旧 §10（四道例题） & 改为独立练习 P1--P4，详解在附录 A；教学示范另换数字，避免提前曝光 \\
旧 §11（易错与速查） & §8，改为八条（新增换基转置一条），并与正文交叉引用 \\
旧 §12（缺口清单） & 本附录 C.3 \\
旧附录 A（指标速查） & 附录 B \\
旧附录 B（自测六问） & 独立练习 P5--P10，答案移入附录 A \\
\end{tabular}
\tabbot
\end{center}

\subhdr{C.3\quad 本稿补充的缺口}

以下每条写四项：\textbf{类型；来源页与原短语；缺了什么；补在哪里}。来源指课件 \texttt{TM2026-Topic8-P1.pdf}。

\gapitem{机制缺口；p9 “Tensor product (Dyad)”}{课件直接给出并矢与它的分量表，没有说明这个对象为什么存在、为什么正好九个分量。补在 §2.2：用双线性定义 + 基并矢展开 + 作用在基矢量对上验证唯一性。}

\gapitem{压缩缺口；p9 “Tensor as a quadratic form: $E=\vect a\ten T\vect a$”}{只给式子，没有说清为什么只有对称部分起作用、为什么等值面是椭球、主轴从哪来。补在 §1.2（对称化证明）与 §1.3（等值面 + 图 1）。}

\gapitem{适用区间；p9 同一行}{默认 $\ten T$ 对称、$\vect a$ 实；非对称时 $E$ 只读出对称部分，这一点课件没写。补在 §1.2 的说明；§4.2、§6.2 各给一次惯量与应力对称的理由。}

\gapitem{机制缺口；p10--p11 “$S(t)^{\mathsf T}S(t)=\sum_k\vect n_k\otimes\vect n_k=\unitd$”}{只给结果，没有展示两个并矢点乘怎么缩并。补在 §3.3：用 §2.8 的并矢点乘规则逐项算出 $S^{\mathsf T}S$ 与 $SS^{\mathsf T}$。}

\gapitem{拼接缺口；p8 的对偶张量与 p12 的角速度}{$\vect a^{\times}$ 在 p8，角速度在 p12，中间没有一处把它们接起来。补在 §2.6（叉乘 = 反对称并矢）与 §3.5（$\boldsymbol\Omega^{\times}=(\dd S/\dd t)S^{\mathsf T}$），并说明“角速度是矢量”来自反对称张量只有三个独立分量。}

\gapitem{机制缺口；p10--p12 输运定理}{课件没有给固定系与体坐标系的导数关系。补在 §5.1：矢量版与并矢版都推，并用 $\ten D=\unitd$ 与 $\ten D=\vect a\otimes\vect b$ 两个核对。}

\gapitem{图缺/图替；p8--p12}{课件全程用符号，没有图：二次型等值面、并矢作用、投影并矢、旋转算子、固定系/体系、惯量主轴、对称陀螺、斜截面分解都靠读者脑补。本稿自绘 8 张图（图 1--图 8），每张都标出与当前公式逐项对应的对象。}

\gapitem{自造缺口（本稿自查）}{本稿新增的解释也检查过：§0.4 的换基公式由三步推导给出，不用 $45^\circ$ 数字例子代替证明；§1.5 的谱定理给出完整陈述，证明标为选读但不影响使用；§4.4 的平行轴定理给出展开推导，不只给结果。}

\subhdr{C.4\quad 选读清单与“跳过检查”}

本稿的选读只有四处：
\begin{enumerate}[label=(\arabic*),leftmargin=1.6em]
\item §1.5 末尾的谱定理 (i)+(iii) 拉格朗日乘子证明；
\item §2.10 末尾的上下标/协变逆变说明；
\item §4.4 末尾的椭球惯量张量；
\item §6.4 末尾的高斯定理说明。
\end{enumerate}
\textbf{跳过检查：}遮住这四处，主线仍然读得通——
§1.5 的定理陈述与“怎么用”在选读之前；§1.6 只用定理的三条结论。
§2.10 的字典在选读之前；§3 之后只用下标。
§4.4 的长方体、薄板、平行轴都在椭球选读之前；§4.5 不用椭球结果。
§6.4 的平衡方程在选读之前；§7 的练习不依赖高斯定理的证明。
四处都在末尾写了“跳过后从哪里继续”，没有把后文要用的定义藏在选读里。

\callout{\textbf{交付说明。}本稿的数值结果、矩阵乘法、特征值、图坐标与部分恒等式由随附的 \texttt{verify.py} 复核，输出见 \texttt{verification.txt}；编译日志见 \texttt{handout.log}。\textbf{讲义交付不等于学习者已掌握，也不等于外部审核已通过}；页面目视验收与学习者试读反馈是另外两件事，见交付报告。}

\flowerrule

{\small\color{inkgray}本稿为 2026-11-11 独立返工稿。主题范围保留，章节、例子、图与表述按本稿重新组织；旧稿的笔误与核对证据记录在 \texttt{source-map.json} 与 \texttt{verification.txt}。}

\end{document}
